% =====================================================================
% 3-р бүлэг  —  Системийн Архитектур ба Зохиомж
% design.tex  |  % =====================================================================
% 3-р бүлэг  —  Системийн Архитектур ба Зохиомж
% design.tex  |  % =====================================================================
% 3-р бүлэг  —  Системийн Архитектур ба Зохиомж
% design.tex  |  % =====================================================================
% 3-р бүлэг  —  Системийн Архитектур ба Зохиомж
% design.tex  |  \input{design} гэж үндсэн файлд оруулна
% Шаардлагатай package-ууд: booktabs, float, graphicx, tabularx,
%   listings, amsmath, tikz (диаграмд), hyperref
% =====================================================================

\chapter{Системийн Архитектур ба Зохиомж}

Системийн архитектурын загварчлал нь зөвхөн технологийн сонголтын асуудал
биш~--- энэ нь байгууллагын хэрэгцээ, бизнесийн хил хязгаар, команд
бүтэц, ирээдүйн тэлэлтийн зорилтыг нийлүүлж, тогтвортой үндсэн шийдвэрт
хувиргах зохион байгуулалтын үйл явц юм. Энэхүү бүлэгт ``юуг яагаад
ингэж бүтэцлэх болсон'' гэсэн асуултад хариулж, системийн гурван үндсэн
давхарга~--- Frontend Microfrontend топологи, Backend Микросервисийн
топологи, Өгөгдлийн давхарга~--- тус бүрийн загварчлалын зарчим, хил
хязгаар, нэгдсэн дизайн шийдвэрийг дэлгэрэнгүй авч үзнэ.

% ─────────────────────────────────────────────────────────────────────
\section{Архитектурын Загварчлалын Үндэслэл}
% ─────────────────────────────────────────────────────────────────────

\subsection{Асуудлын Тодорхойлолт ба Архитектурын Шаардлага}

Монгол Улсын Их Сургуулийн дипломын ажлын удирдлагын одоогийн тогтолцоо
нь Jakarta Server Faces~2.2 (JSF) дээр суурилсан монолит вэб
программ хангамж бөгөөд бизнесийн логик, хуудасны загвар, мэдээллийн
давхарга нь нэгдсэн байдлаар хэрэгжсэн. Энэхүү монолит бүтэц нь системийн
хөгжлийн өмнө дараах тулгамдсан асуудлуудыг үүсгэж байв:

\begin{enumerate}
  \item \textbf{Масштабдах чадварын хязгаарлалт:} JVM heap дээрх
    HttpSession суурилсан тогтолцоо нь зэрэгцэх хэрэглэгчийн тоо
    нэмэгдэх тусам санах ойн хэрэглээ шугаман өсдөг тул өндөр
    ачааллын орчинд нарийн тохируулга шаарддаг.
  \item \textbf{Хөгжүүлэлтийн бие даасан байдлын дутагдал:} Нэг
    хуудасны өөрчлөлт нь бүтэн WAR байршуулалтыг шаарддаг тул
    frontend болон backend хөгжүүлэлт нэгэн зэрэг хийгдэх боломжгүй.
  \item \textbf{Орчин үеийн хэрэглэгчийн интерфэйс дутагдалтай байдал:}
    JSF-ийн request-response мөчлөг нь Single Page Application (SPA)-ийн
    responsive, real-time шинж чанарыг хязгаарладаг.
  \item \textbf{Технологийн хуучин өрийн хуримтлал:} JSF Facelets
    template-ийн нарийн шийдэл нь орчин үеийн TypeScript, React
    экосистемийн хөгжүүлэгчдэд суралцахад хэцүү тул элсүүлэлтийн
    зардал өндөр байдаг.
\end{enumerate}

Энэ асуудлуудад хариулах архитектурын үндсэн шаардлагуудыг Quality
Attribute Workshop (QAW, Barbacci et al., 2003) аргачлалаар тодорхойлж,
дараах таван тэргүүлэх шинж чанарыг сонгов: гүйцэтгэл, засвар
үйлчилгээний чадвар, бие даасан байршуулалтын чадвар, аюулгүй байдал,
шилжилтийн эрсдэлийн хамгаалалт.

\subsection{Архитектурын Стратегийн Сонголт: Аажмын Шилжилтийн Хандлага}

Одоо байгаа JSF системийг шилжүүлэхэд гурван стратегийн сонголт боломжтой:

\begin{enumerate}
  \item \textbf{Бүрэн Дахин Бичих (Big Bang Rewrite):} Одоогийн системийг
    бүрэн орхиж, шинийг эхнээс бичих. Онолын хувьд хамгийн цэвэр боловч
    Fowler (2004)-ийн тодорхойлсноор ``хамгийн эрсдэлтэй нэг юм бол
    бүхлийг нэгэн зэрэг дахин бичих'' учраас бизнесийн тасалдлыг дагуулна.
  \item \textbf{Хэв загварын Шинэчлэл (In-place Modernization):} JSF-ийг
    хадгалан, дотор нь орчин үеийн JavaScript фреймворкийг шингээх.
    Энэ нь JSF-ийн request-response мөчлөгийн хязгаарлалтыг арилгахгүй.
  \item \textbf{Strangler Fig загвар суурилсан Аажмын Шилжилт:} JSF-ийн
    монолитийг ``хост'' хэлбэрт хадгалж, шинэ React MFE модулиудыг
    аажмаар хоёр системийг зэрэгцүүлэн ажиллуулснаар нэгтгэх.
    Fowler (2004) болон Newman (2019)-ийн тодорхойлосон энэ загвар нь
    бизнесийн тасралтгүй байдлыг хамгийн сайн хангадаг.
\end{enumerate}

Энэхүү судалгаа нь гуравдахь сонголт~--- Strangler Fig загвар~--- г
сонгосон бөгөөд үндэслэл нь дараах байна: оюутны бүртгэл, дипломын
хамгаалалтын хуваарь, комиссын удирдлага зэрэг бизнесийн нарийн дараалалт
логикийг системийн ажиллагааг зогсоолгүйгээр шилжүүлэх шаардлага байсан.

\begin{table}[H]
    \centering
    \caption{Архитектурын стратегийн харьцуулалт}
    \label{tab:arch_strategy}
    \renewcommand{\arraystretch}{0.7}
    \begin{tabular}{lccc}
        \toprule
        \textbf{Шалгуур} & \textbf{Big Bang} & \textbf{In-place} & \textbf{Strangler Fig} \\
        \midrule
        Бизнесийн эрсдэл         & Өндөр  & Дунд   & Бага   \\
        Шилжилтийн хугацаа       & Богино & Урт    & Аажмын \\
        Технологийн цэвэр байдал & Өндөр  & Бага   & Дунд   \\
        Командын бие даасан байдал & Бага  & Бага   & Өндөр  \\
        Буцаалтын боломж         & Байхгүй & Хэцүү & Бүрэн  \\
        \bottomrule
    \end{tabular}
\end{table}

% ─────────────────────────────────────────────────────────────────────
\section{Frontend Microfrontend Архитектурын Загварчлал}
% ─────────────────────────────────────────────────────────────────────

\subsection{Гибрид Хостын Загварын Онолын Үндэслэл}

Microfrontend (MFE) архитектур нь backend микросервисийн зарчмуудыг
frontend давхаргад хэрэглэх ойлголт бөгөөд Cam Jackson (2019)-ийн
тодорхойлсноор ``Нэг хуудасны апп-ийг бие даасан, тусдаа байршуулагдах
жижиг фронтэнд апп-уудын нэгтгэл'' гэж тодорхойлогдоно.

Энэхүү системд JSF Tomcat нь нэгдсэн \textbf{Host} (Хост) буюу ``хаалт''
үүрэг гүйцэтгэж, React дээр суурилсан бие даасан модулиуд нь
\textbf{Remote} (Алслагдсан) MFE хэлбэрт ажиллана. Runtime Integration
(Ажиллах цагийн нэгтгэл) стратегийг ашигласан нь нэгтгэлийн цагт биш,
харин браузерт ачааллагдах үед модулиуд динамик байдлаар холбогдоно
гэсэн утгатай бөгөөд энэ нь Дэлхийн тэргүүлэх MFE архитектурын загвар
юм (Mezzalira, 2021).

\begin{figure}[H]
    \centering
    \includegraphics[width=14cm]{fig_hybrid_host_architecture.png}
    \caption{Зураг 3.1. Гибрид хостын загварын ерөнхий бүдүүвч. JSF Tomcat хост нь \texttt{/microfrontends/} замаар React Remote MFE-ийг ачааллаж, браузерт Runtime Integration хэрэгжинэ. Эх сурвалж: зохиогчийн судалгаа, 2026.}
    \label{fig:hybrid_host_arch}
\end{figure}

\subsection{Strangler Fig загварын Хэрэглэлт}

Martin Fowler (2004)-ийн Strangler Fig загвар нь Австралийн чулуун
инжирний мод (Ficus macrophylla)-ний органик загвараас санаа авсан:
шинэ ургамал нь хуучин мод дээр ургаж, аажмаар хучиж, хуучин мод
өөрөө унасны дараа бие даан зогсдог. Программ хангамжийн хүрээнд энэ нь:

\begin{enumerate}
  \item \textbf{Identify:} JSF-ийн аль хэсгийг хамгийн түрүүнд
    шилжүүлэх вэ гэдгийг бизнесийн үнэ цэнэ болон техникийн
    хамаарлын дагуу тодорхойлно.
  \item \textbf{Strangle:} Тухайн хэсгийг React MFE хэлбэрт хэрэгжүүлж,
    JSF-ийн хуудасны дотор ``цэнхэр нүх'' (injection point) хэлбэрт
    оруулна.
  \item \textbf{Kill:} Бүх хэсэг MFE болгон шилжсэний дараа JSF-ийн
    харгалзах хуудсыг устган, MFE-г бие даан ажиллуулна.
\end{enumerate}

Дипломын ажлын системд энэ гурван үе шат нь дараах дарааллаар хэрэгжсэн:
эхлээд нэвтрэлтийн хэлбэр (\texttt{module-auth}), дараа нь оюутны
ажлын хэсэг (\texttt{module-student}), дараа нь багшийн удирдлагын
хэсэг (\texttt{module-teacher}), эцэст нь хорооны хяналт
(\texttt{module-committee}).

\begin{figure}[H]
    \centering
    \includegraphics[width=14cm]{fig_strangler_fig_phases.png}
    \caption{Зураг 3.2. Strangler Fig загварын гурван үе шатны диаграм. Цагаан хэсэг нь JSF монолит, цэнхэр хэсэг нь React MFE, сумнууд нь шилжилтийн чиглэлийг заана. Эх сурвалж: зохиогчийн судалгаа, 2026.}
    \label{fig:strangler_phases}
\end{figure}

\subsection{MFE Модулиудын Загварчлал ба Хариуцлагын Хуваарилалт}

Системийн MFE бүтэц нь ``Bounded Context'' (Хязгаарлагдсан контекст,
Evans, 2003) зарчмын дагуу хуваагдав. Bounded Context нь домейны
мэдлэгийн хил хязгаарыг тодорхойлж, тус тусын домейны хэл, загвар,
хариуцлагыг тодорхой болгоно. Тус системд зургаан бие даасан MFE
модуль загварчлагдав:

\subsubsection{host-jsf (JSF Хост)}

JSF Tomcat сервер нь системийн нэгдсэн хаалт буюу \textbf{Application
Shell} үүрэг гүйцэтгэнэ. Үндсэн зорилго нь URL маршрутлалт, нэвтрэлтийн
хяналт, нийтлэг хуудасны бүтэц (header, navigation, footer) хадгалах,
мөн Remote MFE-ийн \texttt{remoteEntry.js} файлыг динамикаар ачаалах
injection точкийг хангах явдал юм. JSF хостын ``нүх'' бүр нь
\texttt{<div id=\texttt{"}microfrontend-root\texttt{"}/>} хэлбэртэй
DOM цэгт React-ийн render tree суурилна.

\subsubsection{module-auth (Нэвтрэлтийн MFE)}

Нэвтрэх хэлбэр, бүртгэлийн хэлбэр, нууц үг сэргээх урсгал зэргийг
хариуцна. Энэ модуль нь JWT токен авч, \texttt{localStorage}-д хадгалах
хариуцлага хүлээна. Бусад MFE-тэй харьцуулахад \texttt{module-auth} нь
бусад модулиудаас \textbf{тусгаарлагдсан} (isolated) ажиллах ёстой~---
нэвтрэлт хийгдэхэд хостын JSF хуудас руу чиглүүлэлт хийж, улмаар
хэрэглэгчийн дүрд тохирсон MFE ачааллагдана.

\subsubsection{module-student (Оюутны MFE)}

Оюутны хэрэглэгчийн интерфэйсийн бүрэн хариуцлагыг хүлээнэ: дипломын
сэдэв хүсэх, сэдвийн жагсаалт харах, удирдагч багштайгаа харилцах,
явцын тайлан оруулах, эцсийн бүтээлийн файл хавсаргах. Энэ модуль нь
\texttt{thesis\_service} болон \texttt{workflow\_service}-тэй API-аар
холбогдох бөгөөд real-time мэдэгдэл авах үүднээс
\texttt{notification\_service}-тэй WebSocket холболт тогтооно.

\subsubsection{module-teacher (Багшийн MFE)}

Удирдагч багшийн ажлын хэсгийг бүхэлд нь удирдана: оюутны дипломын
ажлыг зөвшөөрөх эсвэл буцаах, явцын тайлантай ажиллах, хэлэлцүүлгийн
хуваарь тохируулах, нийт ажлынхаа дашбоард харах. Багшийн MFE нь
\texttt{user\_service}-аас оюутны мэдээлэл татаж, \texttt{topic\_service}-аас
нэр дэвшигчийн жагсаалт авч ажиллана.

\subsubsection{module-committee (Хорооны MFE)}

Хамгаалалтын комиссын гишүүдэд зориулсан тусгай интерфэйс. Хорооны
бүрэлдэхүүн харах, хэлэлцүүлгийн хуваарийг баталгаажуулах, дүгнэлт
оруулах, нэгдсэн дүн гаргах үүрэгтэй. Энэ модуль нь ролын хязгаарлалт
хамгийн нарийн~--- зөвхөн \texttt{ROLE\_COMMITTEE\_MEMBER} дүр бүхий
хэрэглэгчид нэвтрэх боломжтой.

\subsubsection{module-thesis (Дипломын Ажлын Нэгдсэн MFE)}

Дипломын ажлын бүрэн амьдралын мөчлөгийн хяналт. Оюутан, багш,
хоёуланд нь харагдах нэгдсэн хэлбэрт харагддаг. Явцын Timeline,
баримт бичгийн хувилбарын түүх, тайлбарын гинж (comment thread)
зэргийг агуулна.

\subsubsection{shared-ui (Хуваалцсан UI Сан)}

Бусад MFE-д дундын хэрэгцээтэй бүрэлдэхүүнүүдийг нэгтгэн хадгалах
модуль. \texttt{PageHeader}, \texttt{StatusBadge}, \texttt{DataTable},
\texttt{PortalTabs} зэрэг бүрэлдэхүүнүүд нь энэ модульд хадгалагдаж,
бусад MFE-д \texttt{shared singleton} хэлбэрт нэгтгэгдэнэ. Ant Design
компонент санг тус модуль дамжуулан \textbf{нэгдсэн singleton} байдлаар
хуваалцах нь нийт bundle хэмжээг оновчтой байлгана.

\begin{figure}[H]
    \centering
    \includegraphics[width=14cm]{fig_mfe_module_map.png}
    \caption{Зураг 3.3. MFE модулиудын хариуцлагын хуваарилалт ба хоорондын харилцааны диаграм. Сум бүр нь ``хэрэглэнэ'' (consumes) харилцааг илэрхийлнэ. Эх сурвалж: зохиогчийн судалгаа, 2026.}
    \label{fig:mfe_module_map}
\end{figure}

\subsection{MFE Нэгтгэлийн Стратеги: Runtime Integration}

MFE нэгтгэлийн арга нь гурван ангилалд хуваагдана: \textbf{Build-time
Integration} (нэгтгэлийн үед), \textbf{Server-side Integration} (сервер
талын) болон \textbf{Runtime Integration} (ажиллах үеийн). Тус
системд Runtime Integration сонгосны шалтгаан нь:

\begin{itemize}
  \item Build-time Integration нь бие даасан байршуулалтыг зөвшөөрдөггүй~---
    нэг модуль өөрчлөгдөх бүрт бүгдийг дахин build хийх шаардлагатай болно.
  \item Server-side Integration (SSI) нь JSF-ийн request-response
    мөчлөгтэй зөрчилдөж, template бүрдэлтэд нарийн тохиргоо шаардана.
  \item Runtime Integration нь \textbf{Vite Module Federation} хэрэгслийн
    ачаар браузерт \texttt{remoteEntry.js} динамикаар татаж, зөвхөн
    шаардлагатай модулийг ачаалах ``lazy loading'' зарчмаар ажиллах
    боломжийг олгодог.
\end{itemize}

Runtime Integration-ийн загварт \texttt{remoteEntry.js} нь Remote MFE-ийн
``агуулгын хүснэгт'' (manifest) үүрэг гүйцэтгэж, Host нь энэ файлыг
ачаалснаар тухайн MFE-ийн экспортлосон бүрэлдэхүүнүүдийн мэдрэх зам
(resolution path)-ийг олж авдаг.

% ─────────────────────────────────────────────────────────────────────
\section{Backend Микросервисийн Архитектурын Загварчлал}
% ─────────────────────────────────────────────────────────────────────

\subsection{Монолитоос Микросервис Руу Шилжих Зарчим}

Backend-ийн архитектурын загварчлал нь Sam Newman (2019)-ийн
``Building Microservices'' болон Eric Evans (2003)-ийн Domain-Driven
Design зарчмуудад тулгуурлан хийгдсэн. Монолит backend нь ``Сервис бүр
өөрийн бие даасан байршуулалтын нэгж, нэг бизнесийн чадварыг хариуцна''
гэсэн Mikrosервисийн үндсэн зарчмын дагуу хуваагдсан.

Хуваах шалгуур нь дараах байв: нэгдүгээрт,
Хил хязгаарлагдсан контекст (Bounded Context)~--- домейны мэдлэгийн
хил хязгаар; хоёрдугаарт, Бизнесийн ажлын хэмжээний ялгаа~--- зарим
сервис нь өндөр ачааллыг дааж чаддаг байх шаардлагатай бол зарим нь
тийм ч их хэрэгцээгүй; гуравдугаарт, Командын хариуцлага~--- тус
командын тоо болон ур чадварын хуваарилалттай уялдуулах.

\begin{figure}[H]
    \centering
    \includegraphics[width=14cm]{fig_backend_microservices_overview.png}
    \caption{Зураг 3.4. Backend микросервисийн нэгдсэн бүдүүвч. API Gateway нь нэгдсэн нэвтрэх цэг болж, JWT баталгаажуулалт хийн, тохирох upstream сервис рүү дамжуулна. Эх сурвалж: зохиогчийн судалгаа, 2026.}
    \label{fig:backend_overview}
\end{figure}

\subsection{Микросервисүүдийн Тодорхойлолт ба Хариуцлага}

\subsubsection{auth-service (Баталгаажуулалтын Сервис)}

Хэрэглэгчийн нэвтрэлт, JWT токен үүсгэх, нэвтрэлтийн оролдлого
хянах, нууц үг сэргээх дараалал нь энэ сервисийн цөм хариуцлага юм.
Auth-service нь \texttt{num\_auth-main} багцын дор байрладаг бөгөөд
дотроо гурван дэд сервис агуулна: \texttt{auth-service} (нэвтрэлтийн
цөм), \texttt{authorization-service} (эрхийн удирдлага),
\texttt{user\_detail\_register} (хэрэглэгчийн дэлгэрэнгүй бүртгэл).

JWT-ийн загвар нь Stateless буюу ``Claims суурилсан'' загвар бөгөөд
токен нь \texttt{\{sisiId, roles, iat, exp\}} мэдэгдлийг агуулна.
Токений хэрэглэнгийн хугацаа 24 цаг байхаар тохируулагдсан бөгөөд
Refresh Token механизмийг нэмэлтээр дэмжнэ.

\subsubsection{user\_service (Хэрэглэгчийн Сервис)}

Хэрэглэгчийн профайлын бүрэн мэдрэх хариуцлагыг хүлээнэ: оюутан,
багш, хорооны гишүүн, хэлтсийн мэдээлэл. Auth-service нь зөвхөн
нэвтрэлтийн мэдлэгтэй (\texttt{sisiId}, нууц үг) харьцуулбал
user\_service нь дэлгэрэнгүй профайлын \textbf{единственный эх сурвалж}
(Single Source of Truth) болно. Сервисийн endpoints нь
\texttt{/api/users/students}, \texttt{/api/users/teachers},
\texttt{/api/users/by-email} зэрэг болно.

\subsubsection{topic\_service (Сэдвийн Сервис)}

Дипломын ажлын сэдвийн санд хариуцлагатай: сэдэв зарлах, хүсэлт
хүлээн авах, оюутан-багш хосолт батлах. Энэ сервис нь
\texttt{thesis\_service}-тэй нягт холбоотой боловч тусгаарлагдсан
байдаг~--- сэдвийн амьдралын мөчлөг нь дипломын ажлын амьдралын
мөчлөгтэй ялгаатай.

\subsubsection{thesis\_service (Дипломын Ажлын Сервис)}

Дипломын ажлын үндсэн мета мэдээллийг хадгалах, хувилбар удирдах,
баримт бичгийн хавсаргалт зохицуулах. Дипломын ажлын нэр, хаяг, хар
тэмдэглэл, файлын байршлын мета дата энд хадгалагдана. Файлын биет
агуулга нь object storage (эсвэл локал файлын систем)-д хадгалагдаж,
thesis\_service нь зөвхөн pointer хадгалдаг.

\subsubsection{workflow\_service (Явцын Удирдлагын Сервис)}

Дипломын ажлын амьдралын мөчлөгийн статус шилжилтийг удирдана.
State Machine загварт суурилсан бөгөөд төлөвүүд нь:
\texttt{DRAFT → SUBMITTED → UNDER\_REVIEW → REVISION\_REQUESTED →
APPROVED → DEFENSE\_SCHEDULED → DEFENDED → ARCHIVED}.
Шилжилт бүр нь хугацааны тэмдэглэл, хариуцлагатай хэрэглэгч,
тайлбарын бичлэгийг агуулна.

\subsubsection{evaluation\_service (Үнэлгээний Сервис)}

Комиссын гишүүн бүрийн дүгнэлтийг хүлээн авч, нэгтгэлийн тооцоолол
хийн, эцсийн дүн гаргах. Дүн гаргах алгоритм нь тохируулгатай
(configurable) байхаар загварчлагдсан~--- жишилтийн жин, нийт онооны
хязгаар гэх мэт.

\subsubsection{committee\_service (Комиссын Удирдлагын Сервис)}

Хэлэлцүүлгийн комиссын бүрэлдэхүүн тогтоох, хяналтын хуваарь
гаргах, комиссын гишүүдэд мэдэгдэл илгээх. Энэ сервис нь
\texttt{user\_service}-аас гишүүдийн мэдээлэл авч,
\texttt{notification\_service}-тэй хамтарч мэдэгдэл явуулна.

\subsubsection{API Gateway}

Нэгдсэн нэвтрэх цэг (Single Entry Point) үүрэг гүйцэтгэж, гадаад
хүсэлт бүрийг хүлээн авч, JWT шалгаж, тохирох upstream сервис рүү
дамжуулна (proxy/reverse proxy). Rate limiting, request logging,
CORS хяналт нь gateway давхаргад хэрэгжинэ тул backend сервис тус
бүр энэ логикийг давтан хэрэгжүүлэх шаардлагагүй болно.

\begin{table}[H]
    \centering
    \caption{Backend микросервисүүдийн порт болон хариуцлагын хуваарь}
    \label{tab:service_ports}
    \renewcommand{\arraystretch}{0.7}
    \begin{tabular}{llll}
        \toprule
        \textbf{Сервис} & \textbf{Порт} & \textbf{Технологи} & \textbf{Гол хариуцлага} \\
        \midrule
        auth-service         & 8887 & Spring Boot  & JWT үүсгэх, нэвтрэлт \\
        user\_service        & 8086 & Spring WebFlux & Хэрэглэгчийн профайл \\
        topic\_service       & 8081 & Spring WebFlux & Сэдвийн удирдлага \\
        thesis\_service      & 8083 & Spring WebFlux & Дипломын ажлын мета \\
        workflow\_service    & 8084 & Spring WebFlux & Статус шилжилт \\
        evaluation\_service  & 8085 & Spring WebFlux & Дүгнэлт, үнэлгээ \\
        committee\_service   & 8082 & Spring WebFlux & Комиссын удирдлага \\
        Tomcat (JSF+MFE)     & 8080 & Apache Tomcat  & Host, static file serve \\
        \bottomrule
    \end{tabular}
\end{table}

\subsection{Сервисүүдийн Хоорондын Харилцааны Загвар}

Микросервисүүдийн хоорондын харилцааны загварчлалд хоёр үндсэн хэв
загвар хэрэглэгдэв:

\textbf{Синхрон харилцаа (Synchronous Communication):} HTTP/REST
протоколоор шууд дуудлага. \texttt{workflow\_service} нь статус
шилжилт хийхдээ \texttt{thesis\_service}-г шууд дуудаж дипломын
ажлын бичлэгийг шинэчлэнэ. Ийм шууд харилцааны тохиолдолд Circuit
Breaker загварыг (Nygard, 2007) хэрэглэж, upstream сервисийн
доройтол нь каскад байдлаар бусдад нэвтрэхээс хамгаалдаг.

\textbf{Асинхрон харилцаа (Asynchronous Communication):} Мэдэгдэл
(notification) явуулах, аудит бичлэг хийх гэх мэт цаг хугацааны
хамаарал бага үйлдлүүд нь Event-Driven загвараар хэрэгжих ёстой.
Ирээдүйн хөгжлийн шатанд Apache Kafka Message Broker нэвтрүүлснээр
\texttt{workflow\_service}-ийн статус шилжилтийн event нь
\texttt{notification\_service}-т асинхрон байдлаар хүрэх болно.

\begin{figure}[H]
    \centering
    \includegraphics[width=14cm]{fig_service_communication.png}
    \caption{Зураг 3.5. Микросервисүүдийн харилцааны диаграм. Хатуу шугам нь синхрон REST дуудлага, тасархай шугам нь асинхрон event-г тэмдэглэнэ. Эх сурвалж: зохиогчийн судалгаа, 2026.}
    \label{fig:service_communication}
\end{figure}

% ─────────────────────────────────────────────────────────────────────
\section{Өгөгдлийн Давхаргын Загварчлал}
% ─────────────────────────────────────────────────────────────────────

\subsection{Database Per Service загварын Сонголт}

Микросервисийн архитектурт өгөгдлийн давхаргыг хэрхэн хуваах вэ
гэсэн асуудал нь хамгийн нарийн төвөгтэй архитектурын шийдвэрүүдийн
нэг юм. Хоёр үндсэн сонголт байдаг:

\begin{enumerate}
  \item \textbf{Хуваалцсан өгөгдлийн сан (Shared Database):} Бүх
    сервис нэг өгөгдлийн санг хуваалцна. Хэрэгжүүлэлт энгийн боловч
    сервисүүдийн хооронд нягт холбоос (tight coupling) үүсдэг~---
    нэг сервисийн schema өөрчлөлт нь бусдад нөлөөлнө.
  \item \textbf{Сервис тус бүрийн өгөгдлийн сан (Database Per Service):}
    Сервис бүр өөрийн тусдаа schema эзэмшинэ. Сервис бүр бие даан
    deploy хийгдэх, технологийн сонголтоо чөлөөтэй хийх боломжтой
    боловч Cross-service join, distributed transaction нарийн болно.
\end{enumerate}

Энэхүү судалгаанд прагматик шийдлийг сонгосон: бүх сервис нэг
PostgreSQL instance-д ажилладаг боловч \textbf{schema-level тусгаарлалт}
хийгдсэн~--- сервис бүр өөрийн тусдаа schema-д бичиж, cross-schema
foreign key ашиглахгүй. Энэ нь ``Database Per Service'' загварын бие
даасан байдлын ашиг тусыг хэсэгчлэн хадгалахын зэрэгцээ нэг instance-ийн
хялбар засвар үйлчилгээний давуу талыг авч үлдэнэ.

\begin{figure}[H]
    \centering
    \includegraphics[width=14cm]{fig_database_schema_isolation.png}
    \caption{Зураг 3.6. Schema-level тусгаарлалтын загвар. Нэг PostgreSQL instance дотор сервис бүр өөрийн schema эзэмших бөгөөд хоорондын хамаарал нь application давхаргад шийдэгдэнэ. Эх сурвалж: зохиогчийн судалгаа, 2026.}
    \label{fig:db_schema_isolation}
\end{figure}

\subsection{Домейн Загварын Тодорхойлолт}

Системийн гол домейн объектуудын хоорондын харилцааг Entity-Relationship
(ER) загварт тодорхойлов. Хэдийгээр Cross-service join ашиглахгүй ч,
бизнесийн логикийн хэрэгцээгээр тодорхой домейн хооронд нягт харилцаа
оршино. Энэ харилцааг application давхаргад API дуудлагаар шийдвэрлэнэ.

Гол домейн объектууд болон тэдгээрийн атрибутуудыг дараах байдлаар
тодорхойлов:

\begin{itemize}
  \item \textbf{User} (user\_service): \texttt{id, sisiId, firstName,
    lastName, email, role, departmentId}
  \item \textbf{Department} (user\_service): \texttt{id, name, code,
    facultyId}
  \item \textbf{Topic} (topic\_service): \texttt{id, title, description,
    teacherId, maxStudents, status, academicYear}
  \item \textbf{Thesis} (thesis\_service): \texttt{id, title, studentId,
    teacherId, topicId, fileUrl, version, submittedAt}
  \item \textbf{WorkflowState} (workflow\_service): \texttt{id, thesisId,
    fromStatus, toStatus, actorId, comment, transitionedAt}
  \item \textbf{Committee} (committee\_service): \texttt{id, name,
    defenseDate, venueId, chairId}
  \item \textbf{CommitteeMember} (committee\_service): \texttt{id,
    committeeId, userId, role}
  \item \textbf{Evaluation} (evaluation\_service): \texttt{id, thesisId,
    committeeId, evaluatorId, score, criteria, submittedAt}
\end{itemize}

\begin{figure}[H]
    \centering
    \includegraphics[width=14cm]{fig_domain_er_diagram.png}
    \caption{Зураг 3.7. Системийн домейн загварын ER диаграм. Хатуу шугам нь нэг schema дотрох foreign key, тасархай шугам нь application давхаргын logical харилцааг заана. Эх сурвалж: зохиогчийн судалгаа, 2026.}
    \label{fig:domain_er}
\end{figure}

\subsection{Мэдээллийн Нийцэл ба Eventual Consistency}

Cross-service мэдээллийн нийцэлд ``Eventual Consistency'' (Аажмын
нийцэл) зарчмыг хэрэглэнэ. Distributed Transaction (2-Phase Commit)
нь микросервисийн орчинд хэдийгээр боломжтой боловч хүнд хэцүү
тогтолцоо шаарддаг тул Saga Pattern (Garcia-Molina \& Salem, 1987)
болон Outbox Pattern-ийг илүүд үзнэ.

Жишээ нь: оюутан дипломын ажлаа \texttt{SUBMITTED} төлөвт шилжүүлэхэд
\texttt{workflow\_service} нь шилжилтийн бичлэгийг хадгалж, дараа нь
\texttt{thesis\_service}-г шинэчлэх хүсэлт илгээнэ. Хэрэв
\texttt{thesis\_service} хариу өгөхгүй бол \texttt{workflow\_service}
нь retry логикоор дахин оролдох бөгөөд эцэст нь нийцэл тогтоно~---
энэ нь ``Eventual'' буюу аажмын нийцэл юм.

% ─────────────────────────────────────────────────────────────────────
\section{Нэгдсэн Архитектурын Загварын Нэгтгэл}
% ─────────────────────────────────────────────────────────────────────

Дээр тодорхойлсон гурван давхаргын загварчлал нэгдсэн байдлаар ажиллахад
дараах нэгдсэн зарчмуудыг дагана:

\begin{enumerate}
  \item \textbf{Нэг хариуцлага:} MFE модуль болон backend сервис тус
    бүр нэг бизнесийн чадварыг хариуцна~--- хариуцлагыг ``кэйс''
    дагуу биш, ``домейн'' дагуу хуваана.
  \item \textbf{API гэрээгээр харилцах:} Сервис болон MFE хоорондын
    харилцааны цорын ганц хэлбэр нь баримтжуулсан API гэрээ (OpenAPI
    Specification) юм~--- шууд function call, shared memory гэж байхгүй.
  \item \textbf{Бие даасан байршуулалт:} Сервис болон MFE тус бүр
    бусдыг зогсоолгүй шинэчлэгдэх чадвартай байна.
  \item \textbf{Барилгын бие даасан байдал:} Хөгжүүлэгч баг тус бүр
    өөрийн домейнийг бие даан эзэмшиж, технологийн сонголт дотор
    ухаалаг шийдэл гаргах эрхтэй.
\end{enumerate}

\begin{figure}[H]
    \centering
    \includegraphics[width=14cm]{fig_full_architecture_overview.png}
    \caption{Зураг 3.8. Системийн бүрэн архитектурын нэгдсэн бүдүүвч. Гурван давхарга (Frontend MFE, API Gateway, Backend Сервис) ба өгөгдлийн давхарга нэгдсэн байдлаар харагдана. Эх сурвалж: зохиогчийн судалгаа, 2026.}
    \label{fig:full_architecture}
\end{figure}

3-р бүлгийн загварчлалын шийдвэрүүд нь 4-р бүлгийн практик хэрэгжүүлэлтийн
суурь болно. Дараах бүлэгт ``юуг яагаад ийм загварчилсан'' гэдгийг
``хэрхэн бодитоор кодлон хэрэгжүүлсэн'' гэдэгтэй холбон тайлбарлана.
 гэж үндсэн файлд оруулна
% Шаардлагатай package-ууд: booktabs, float, graphicx, tabularx,
%   listings, amsmath, tikz (диаграмд), hyperref
% =====================================================================

\chapter{Системийн Архитектур ба Зохиомж}

Системийн архитектурын загварчлал нь зөвхөн технологийн сонголтын асуудал
биш~--- энэ нь байгууллагын хэрэгцээ, бизнесийн хил хязгаар, команд
бүтэц, ирээдүйн тэлэлтийн зорилтыг нийлүүлж, тогтвортой үндсэн шийдвэрт
хувиргах зохион байгуулалтын үйл явц юм. Энэхүү бүлэгт ``юуг яагаад
ингэж бүтэцлэх болсон'' гэсэн асуултад хариулж, системийн гурван үндсэн
давхарга~--- Frontend Microfrontend топологи, Backend Микросервисийн
топологи, Өгөгдлийн давхарга~--- тус бүрийн загварчлалын зарчим, хил
хязгаар, нэгдсэн дизайн шийдвэрийг дэлгэрэнгүй авч үзнэ.

% ─────────────────────────────────────────────────────────────────────
\section{Архитектурын Загварчлалын Үндэслэл}
% ─────────────────────────────────────────────────────────────────────

\subsection{Асуудлын Тодорхойлолт ба Архитектурын Шаардлага}

Монгол Улсын Их Сургуулийн дипломын ажлын удирдлагын одоогийн тогтолцоо
нь Jakarta Server Faces~2.2 (JSF) дээр суурилсан монолит вэб
программ хангамж бөгөөд бизнесийн логик, хуудасны загвар, мэдээллийн
давхарга нь нэгдсэн байдлаар хэрэгжсэн. Энэхүү монолит бүтэц нь системийн
хөгжлийн өмнө дараах тулгамдсан асуудлуудыг үүсгэж байв:

\begin{enumerate}
  \item \textbf{Масштабдах чадварын хязгаарлалт:} JVM heap дээрх
    HttpSession суурилсан тогтолцоо нь зэрэгцэх хэрэглэгчийн тоо
    нэмэгдэх тусам санах ойн хэрэглээ шугаман өсдөг тул өндөр
    ачааллын орчинд нарийн тохируулга шаарддаг.
  \item \textbf{Хөгжүүлэлтийн бие даасан байдлын дутагдал:} Нэг
    хуудасны өөрчлөлт нь бүтэн WAR байршуулалтыг шаарддаг тул
    frontend болон backend хөгжүүлэлт нэгэн зэрэг хийгдэх боломжгүй.
  \item \textbf{Орчин үеийн хэрэглэгчийн интерфэйс дутагдалтай байдал:}
    JSF-ийн request-response мөчлөг нь Single Page Application (SPA)-ийн
    responsive, real-time шинж чанарыг хязгаарладаг.
  \item \textbf{Технологийн хуучин өрийн хуримтлал:} JSF Facelets
    template-ийн нарийн шийдэл нь орчин үеийн TypeScript, React
    экосистемийн хөгжүүлэгчдэд суралцахад хэцүү тул элсүүлэлтийн
    зардал өндөр байдаг.
\end{enumerate}

Энэ асуудлуудад хариулах архитектурын үндсэн шаардлагуудыг Quality
Attribute Workshop (QAW, Barbacci et al., 2003) аргачлалаар тодорхойлж,
дараах таван тэргүүлэх шинж чанарыг сонгов: гүйцэтгэл, засвар
үйлчилгээний чадвар, бие даасан байршуулалтын чадвар, аюулгүй байдал,
шилжилтийн эрсдэлийн хамгаалалт.

\subsection{Архитектурын Стратегийн Сонголт: Аажмын Шилжилтийн Хандлага}

Одоо байгаа JSF системийг шилжүүлэхэд гурван стратегийн сонголт боломжтой:

\begin{enumerate}
  \item \textbf{Бүрэн Дахин Бичих (Big Bang Rewrite):} Одоогийн системийг
    бүрэн орхиж, шинийг эхнээс бичих. Онолын хувьд хамгийн цэвэр боловч
    Fowler (2004)-ийн тодорхойлсноор ``хамгийн эрсдэлтэй нэг юм бол
    бүхлийг нэгэн зэрэг дахин бичих'' учраас бизнесийн тасалдлыг дагуулна.
  \item \textbf{Хэв загварын Шинэчлэл (In-place Modernization):} JSF-ийг
    хадгалан, дотор нь орчин үеийн JavaScript фреймворкийг шингээх.
    Энэ нь JSF-ийн request-response мөчлөгийн хязгаарлалтыг арилгахгүй.
  \item \textbf{Strangler Fig загвар суурилсан Аажмын Шилжилт:} JSF-ийн
    монолитийг ``хост'' хэлбэрт хадгалж, шинэ React MFE модулиудыг
    аажмаар хоёр системийг зэрэгцүүлэн ажиллуулснаар нэгтгэх.
    Fowler (2004) болон Newman (2019)-ийн тодорхойлосон энэ загвар нь
    бизнесийн тасралтгүй байдлыг хамгийн сайн хангадаг.
\end{enumerate}

Энэхүү судалгаа нь гуравдахь сонголт~--- Strangler Fig загвар~--- г
сонгосон бөгөөд үндэслэл нь дараах байна: оюутны бүртгэл, дипломын
хамгаалалтын хуваарь, комиссын удирдлага зэрэг бизнесийн нарийн дараалалт
логикийг системийн ажиллагааг зогсоолгүйгээр шилжүүлэх шаардлага байсан.

\begin{table}[H]
    \centering
    \caption{Архитектурын стратегийн харьцуулалт}
    \label{tab:arch_strategy}
    \renewcommand{\arraystretch}{0.7}
    \begin{tabular}{lccc}
        \toprule
        \textbf{Шалгуур} & \textbf{Big Bang} & \textbf{In-place} & \textbf{Strangler Fig} \\
        \midrule
        Бизнесийн эрсдэл         & Өндөр  & Дунд   & Бага   \\
        Шилжилтийн хугацаа       & Богино & Урт    & Аажмын \\
        Технологийн цэвэр байдал & Өндөр  & Бага   & Дунд   \\
        Командын бие даасан байдал & Бага  & Бага   & Өндөр  \\
        Буцаалтын боломж         & Байхгүй & Хэцүү & Бүрэн  \\
        \bottomrule
    \end{tabular}
\end{table}

% ─────────────────────────────────────────────────────────────────────
\section{Frontend Microfrontend Архитектурын Загварчлал}
% ─────────────────────────────────────────────────────────────────────

\subsection{Гибрид Хостын Загварын Онолын Үндэслэл}

Microfrontend (MFE) архитектур нь backend микросервисийн зарчмуудыг
frontend давхаргад хэрэглэх ойлголт бөгөөд Cam Jackson (2019)-ийн
тодорхойлсноор ``Нэг хуудасны апп-ийг бие даасан, тусдаа байршуулагдах
жижиг фронтэнд апп-уудын нэгтгэл'' гэж тодорхойлогдоно.

Энэхүү системд JSF Tomcat нь нэгдсэн \textbf{Host} (Хост) буюу ``хаалт''
үүрэг гүйцэтгэж, React дээр суурилсан бие даасан модулиуд нь
\textbf{Remote} (Алслагдсан) MFE хэлбэрт ажиллана. Runtime Integration
(Ажиллах цагийн нэгтгэл) стратегийг ашигласан нь нэгтгэлийн цагт биш,
харин браузерт ачааллагдах үед модулиуд динамик байдлаар холбогдоно
гэсэн утгатай бөгөөд энэ нь Дэлхийн тэргүүлэх MFE архитектурын загвар
юм (Mezzalira, 2021).

\begin{figure}[H]
    \centering
    \includegraphics[width=14cm]{fig_hybrid_host_architecture.png}
    \caption{Зураг 3.1. Гибрид хостын загварын ерөнхий бүдүүвч. JSF Tomcat хост нь \texttt{/microfrontends/} замаар React Remote MFE-ийг ачааллаж, браузерт Runtime Integration хэрэгжинэ. Эх сурвалж: зохиогчийн судалгаа, 2026.}
    \label{fig:hybrid_host_arch}
\end{figure}

\subsection{Strangler Fig загварын Хэрэглэлт}

Martin Fowler (2004)-ийн Strangler Fig загвар нь Австралийн чулуун
инжирний мод (Ficus macrophylla)-ний органик загвараас санаа авсан:
шинэ ургамал нь хуучин мод дээр ургаж, аажмаар хучиж, хуучин мод
өөрөө унасны дараа бие даан зогсдог. Программ хангамжийн хүрээнд энэ нь:

\begin{enumerate}
  \item \textbf{Identify:} JSF-ийн аль хэсгийг хамгийн түрүүнд
    шилжүүлэх вэ гэдгийг бизнесийн үнэ цэнэ болон техникийн
    хамаарлын дагуу тодорхойлно.
  \item \textbf{Strangle:} Тухайн хэсгийг React MFE хэлбэрт хэрэгжүүлж,
    JSF-ийн хуудасны дотор ``цэнхэр нүх'' (injection point) хэлбэрт
    оруулна.
  \item \textbf{Kill:} Бүх хэсэг MFE болгон шилжсэний дараа JSF-ийн
    харгалзах хуудсыг устган, MFE-г бие даан ажиллуулна.
\end{enumerate}

Дипломын ажлын системд энэ гурван үе шат нь дараах дарааллаар хэрэгжсэн:
эхлээд нэвтрэлтийн хэлбэр (\texttt{module-auth}), дараа нь оюутны
ажлын хэсэг (\texttt{module-student}), дараа нь багшийн удирдлагын
хэсэг (\texttt{module-teacher}), эцэст нь хорооны хяналт
(\texttt{module-committee}).

\begin{figure}[H]
    \centering
    \includegraphics[width=14cm]{fig_strangler_fig_phases.png}
    \caption{Зураг 3.2. Strangler Fig загварын гурван үе шатны диаграм. Цагаан хэсэг нь JSF монолит, цэнхэр хэсэг нь React MFE, сумнууд нь шилжилтийн чиглэлийг заана. Эх сурвалж: зохиогчийн судалгаа, 2026.}
    \label{fig:strangler_phases}
\end{figure}

\subsection{MFE Модулиудын Загварчлал ба Хариуцлагын Хуваарилалт}

Системийн MFE бүтэц нь ``Bounded Context'' (Хязгаарлагдсан контекст,
Evans, 2003) зарчмын дагуу хуваагдав. Bounded Context нь домейны
мэдлэгийн хил хязгаарыг тодорхойлж, тус тусын домейны хэл, загвар,
хариуцлагыг тодорхой болгоно. Тус системд зургаан бие даасан MFE
модуль загварчлагдав:

\subsubsection{host-jsf (JSF Хост)}

JSF Tomcat сервер нь системийн нэгдсэн хаалт буюу \textbf{Application
Shell} үүрэг гүйцэтгэнэ. Үндсэн зорилго нь URL маршрутлалт, нэвтрэлтийн
хяналт, нийтлэг хуудасны бүтэц (header, navigation, footer) хадгалах,
мөн Remote MFE-ийн \texttt{remoteEntry.js} файлыг динамикаар ачаалах
injection точкийг хангах явдал юм. JSF хостын ``нүх'' бүр нь
\texttt{<div id=\texttt{"}microfrontend-root\texttt{"}/>} хэлбэртэй
DOM цэгт React-ийн render tree суурилна.

\subsubsection{module-auth (Нэвтрэлтийн MFE)}

Нэвтрэх хэлбэр, бүртгэлийн хэлбэр, нууц үг сэргээх урсгал зэргийг
хариуцна. Энэ модуль нь JWT токен авч, \texttt{localStorage}-д хадгалах
хариуцлага хүлээна. Бусад MFE-тэй харьцуулахад \texttt{module-auth} нь
бусад модулиудаас \textbf{тусгаарлагдсан} (isolated) ажиллах ёстой~---
нэвтрэлт хийгдэхэд хостын JSF хуудас руу чиглүүлэлт хийж, улмаар
хэрэглэгчийн дүрд тохирсон MFE ачааллагдана.

\subsubsection{module-student (Оюутны MFE)}

Оюутны хэрэглэгчийн интерфэйсийн бүрэн хариуцлагыг хүлээнэ: дипломын
сэдэв хүсэх, сэдвийн жагсаалт харах, удирдагч багштайгаа харилцах,
явцын тайлан оруулах, эцсийн бүтээлийн файл хавсаргах. Энэ модуль нь
\texttt{thesis\_service} болон \texttt{workflow\_service}-тэй API-аар
холбогдох бөгөөд real-time мэдэгдэл авах үүднээс
\texttt{notification\_service}-тэй WebSocket холболт тогтооно.

\subsubsection{module-teacher (Багшийн MFE)}

Удирдагч багшийн ажлын хэсгийг бүхэлд нь удирдана: оюутны дипломын
ажлыг зөвшөөрөх эсвэл буцаах, явцын тайлантай ажиллах, хэлэлцүүлгийн
хуваарь тохируулах, нийт ажлынхаа дашбоард харах. Багшийн MFE нь
\texttt{user\_service}-аас оюутны мэдээлэл татаж, \texttt{topic\_service}-аас
нэр дэвшигчийн жагсаалт авч ажиллана.

\subsubsection{module-committee (Хорооны MFE)}

Хамгаалалтын комиссын гишүүдэд зориулсан тусгай интерфэйс. Хорооны
бүрэлдэхүүн харах, хэлэлцүүлгийн хуваарийг баталгаажуулах, дүгнэлт
оруулах, нэгдсэн дүн гаргах үүрэгтэй. Энэ модуль нь ролын хязгаарлалт
хамгийн нарийн~--- зөвхөн \texttt{ROLE\_COMMITTEE\_MEMBER} дүр бүхий
хэрэглэгчид нэвтрэх боломжтой.

\subsubsection{module-thesis (Дипломын Ажлын Нэгдсэн MFE)}

Дипломын ажлын бүрэн амьдралын мөчлөгийн хяналт. Оюутан, багш,
хоёуланд нь харагдах нэгдсэн хэлбэрт харагддаг. Явцын Timeline,
баримт бичгийн хувилбарын түүх, тайлбарын гинж (comment thread)
зэргийг агуулна.

\subsubsection{shared-ui (Хуваалцсан UI Сан)}

Бусад MFE-д дундын хэрэгцээтэй бүрэлдэхүүнүүдийг нэгтгэн хадгалах
модуль. \texttt{PageHeader}, \texttt{StatusBadge}, \texttt{DataTable},
\texttt{PortalTabs} зэрэг бүрэлдэхүүнүүд нь энэ модульд хадгалагдаж,
бусад MFE-д \texttt{shared singleton} хэлбэрт нэгтгэгдэнэ. Ant Design
компонент санг тус модуль дамжуулан \textbf{нэгдсэн singleton} байдлаар
хуваалцах нь нийт bundle хэмжээг оновчтой байлгана.

\begin{figure}[H]
    \centering
    \includegraphics[width=14cm]{fig_mfe_module_map.png}
    \caption{Зураг 3.3. MFE модулиудын хариуцлагын хуваарилалт ба хоорондын харилцааны диаграм. Сум бүр нь ``хэрэглэнэ'' (consumes) харилцааг илэрхийлнэ. Эх сурвалж: зохиогчийн судалгаа, 2026.}
    \label{fig:mfe_module_map}
\end{figure}

\subsection{MFE Нэгтгэлийн Стратеги: Runtime Integration}

MFE нэгтгэлийн арга нь гурван ангилалд хуваагдана: \textbf{Build-time
Integration} (нэгтгэлийн үед), \textbf{Server-side Integration} (сервер
талын) болон \textbf{Runtime Integration} (ажиллах үеийн). Тус
системд Runtime Integration сонгосны шалтгаан нь:

\begin{itemize}
  \item Build-time Integration нь бие даасан байршуулалтыг зөвшөөрдөггүй~---
    нэг модуль өөрчлөгдөх бүрт бүгдийг дахин build хийх шаардлагатай болно.
  \item Server-side Integration (SSI) нь JSF-ийн request-response
    мөчлөгтэй зөрчилдөж, template бүрдэлтэд нарийн тохиргоо шаардана.
  \item Runtime Integration нь \textbf{Vite Module Federation} хэрэгслийн
    ачаар браузерт \texttt{remoteEntry.js} динамикаар татаж, зөвхөн
    шаардлагатай модулийг ачаалах ``lazy loading'' зарчмаар ажиллах
    боломжийг олгодог.
\end{itemize}

Runtime Integration-ийн загварт \texttt{remoteEntry.js} нь Remote MFE-ийн
``агуулгын хүснэгт'' (manifest) үүрэг гүйцэтгэж, Host нь энэ файлыг
ачаалснаар тухайн MFE-ийн экспортлосон бүрэлдэхүүнүүдийн мэдрэх зам
(resolution path)-ийг олж авдаг.

% ─────────────────────────────────────────────────────────────────────
\section{Backend Микросервисийн Архитектурын Загварчлал}
% ─────────────────────────────────────────────────────────────────────

\subsection{Монолитоос Микросервис Руу Шилжих Зарчим}

Backend-ийн архитектурын загварчлал нь Sam Newman (2019)-ийн
``Building Microservices'' болон Eric Evans (2003)-ийн Domain-Driven
Design зарчмуудад тулгуурлан хийгдсэн. Монолит backend нь ``Сервис бүр
өөрийн бие даасан байршуулалтын нэгж, нэг бизнесийн чадварыг хариуцна''
гэсэн Mikrosервисийн үндсэн зарчмын дагуу хуваагдсан.

Хуваах шалгуур нь дараах байв: нэгдүгээрт,
Хил хязгаарлагдсан контекст (Bounded Context)~--- домейны мэдлэгийн
хил хязгаар; хоёрдугаарт, Бизнесийн ажлын хэмжээний ялгаа~--- зарим
сервис нь өндөр ачааллыг дааж чаддаг байх шаардлагатай бол зарим нь
тийм ч их хэрэгцээгүй; гуравдугаарт, Командын хариуцлага~--- тус
командын тоо болон ур чадварын хуваарилалттай уялдуулах.

\begin{figure}[H]
    \centering
    \includegraphics[width=14cm]{fig_backend_microservices_overview.png}
    \caption{Зураг 3.4. Backend микросервисийн нэгдсэн бүдүүвч. API Gateway нь нэгдсэн нэвтрэх цэг болж, JWT баталгаажуулалт хийн, тохирох upstream сервис рүү дамжуулна. Эх сурвалж: зохиогчийн судалгаа, 2026.}
    \label{fig:backend_overview}
\end{figure}

\subsection{Микросервисүүдийн Тодорхойлолт ба Хариуцлага}

\subsubsection{auth-service (Баталгаажуулалтын Сервис)}

Хэрэглэгчийн нэвтрэлт, JWT токен үүсгэх, нэвтрэлтийн оролдлого
хянах, нууц үг сэргээх дараалал нь энэ сервисийн цөм хариуцлага юм.
Auth-service нь \texttt{num\_auth-main} багцын дор байрладаг бөгөөд
дотроо гурван дэд сервис агуулна: \texttt{auth-service} (нэвтрэлтийн
цөм), \texttt{authorization-service} (эрхийн удирдлага),
\texttt{user\_detail\_register} (хэрэглэгчийн дэлгэрэнгүй бүртгэл).

JWT-ийн загвар нь Stateless буюу ``Claims суурилсан'' загвар бөгөөд
токен нь \texttt{\{sisiId, roles, iat, exp\}} мэдэгдлийг агуулна.
Токений хэрэглэнгийн хугацаа 24 цаг байхаар тохируулагдсан бөгөөд
Refresh Token механизмийг нэмэлтээр дэмжнэ.

\subsubsection{user\_service (Хэрэглэгчийн Сервис)}

Хэрэглэгчийн профайлын бүрэн мэдрэх хариуцлагыг хүлээнэ: оюутан,
багш, хорооны гишүүн, хэлтсийн мэдээлэл. Auth-service нь зөвхөн
нэвтрэлтийн мэдлэгтэй (\texttt{sisiId}, нууц үг) харьцуулбал
user\_service нь дэлгэрэнгүй профайлын \textbf{единственный эх сурвалж}
(Single Source of Truth) болно. Сервисийн endpoints нь
\texttt{/api/users/students}, \texttt{/api/users/teachers},
\texttt{/api/users/by-email} зэрэг болно.

\subsubsection{topic\_service (Сэдвийн Сервис)}

Дипломын ажлын сэдвийн санд хариуцлагатай: сэдэв зарлах, хүсэлт
хүлээн авах, оюутан-багш хосолт батлах. Энэ сервис нь
\texttt{thesis\_service}-тэй нягт холбоотой боловч тусгаарлагдсан
байдаг~--- сэдвийн амьдралын мөчлөг нь дипломын ажлын амьдралын
мөчлөгтэй ялгаатай.

\subsubsection{thesis\_service (Дипломын Ажлын Сервис)}

Дипломын ажлын үндсэн мета мэдээллийг хадгалах, хувилбар удирдах,
баримт бичгийн хавсаргалт зохицуулах. Дипломын ажлын нэр, хаяг, хар
тэмдэглэл, файлын байршлын мета дата энд хадгалагдана. Файлын биет
агуулга нь object storage (эсвэл локал файлын систем)-д хадгалагдаж,
thesis\_service нь зөвхөн pointer хадгалдаг.

\subsubsection{workflow\_service (Явцын Удирдлагын Сервис)}

Дипломын ажлын амьдралын мөчлөгийн статус шилжилтийг удирдана.
State Machine загварт суурилсан бөгөөд төлөвүүд нь:
\texttt{DRAFT → SUBMITTED → UNDER\_REVIEW → REVISION\_REQUESTED →
APPROVED → DEFENSE\_SCHEDULED → DEFENDED → ARCHIVED}.
Шилжилт бүр нь хугацааны тэмдэглэл, хариуцлагатай хэрэглэгч,
тайлбарын бичлэгийг агуулна.

\subsubsection{evaluation\_service (Үнэлгээний Сервис)}

Комиссын гишүүн бүрийн дүгнэлтийг хүлээн авч, нэгтгэлийн тооцоолол
хийн, эцсийн дүн гаргах. Дүн гаргах алгоритм нь тохируулгатай
(configurable) байхаар загварчлагдсан~--- жишилтийн жин, нийт онооны
хязгаар гэх мэт.

\subsubsection{committee\_service (Комиссын Удирдлагын Сервис)}

Хэлэлцүүлгийн комиссын бүрэлдэхүүн тогтоох, хяналтын хуваарь
гаргах, комиссын гишүүдэд мэдэгдэл илгээх. Энэ сервис нь
\texttt{user\_service}-аас гишүүдийн мэдээлэл авч,
\texttt{notification\_service}-тэй хамтарч мэдэгдэл явуулна.

\subsubsection{API Gateway}

Нэгдсэн нэвтрэх цэг (Single Entry Point) үүрэг гүйцэтгэж, гадаад
хүсэлт бүрийг хүлээн авч, JWT шалгаж, тохирох upstream сервис рүү
дамжуулна (proxy/reverse proxy). Rate limiting, request logging,
CORS хяналт нь gateway давхаргад хэрэгжинэ тул backend сервис тус
бүр энэ логикийг давтан хэрэгжүүлэх шаардлагагүй болно.

\begin{table}[H]
    \centering
    \caption{Backend микросервисүүдийн порт болон хариуцлагын хуваарь}
    \label{tab:service_ports}
    \renewcommand{\arraystretch}{0.7}
    \begin{tabular}{llll}
        \toprule
        \textbf{Сервис} & \textbf{Порт} & \textbf{Технологи} & \textbf{Гол хариуцлага} \\
        \midrule
        auth-service         & 8887 & Spring Boot  & JWT үүсгэх, нэвтрэлт \\
        user\_service        & 8086 & Spring WebFlux & Хэрэглэгчийн профайл \\
        topic\_service       & 8081 & Spring WebFlux & Сэдвийн удирдлага \\
        thesis\_service      & 8083 & Spring WebFlux & Дипломын ажлын мета \\
        workflow\_service    & 8084 & Spring WebFlux & Статус шилжилт \\
        evaluation\_service  & 8085 & Spring WebFlux & Дүгнэлт, үнэлгээ \\
        committee\_service   & 8082 & Spring WebFlux & Комиссын удирдлага \\
        Tomcat (JSF+MFE)     & 8080 & Apache Tomcat  & Host, static file serve \\
        \bottomrule
    \end{tabular}
\end{table}

\subsection{Сервисүүдийн Хоорондын Харилцааны Загвар}

Микросервисүүдийн хоорондын харилцааны загварчлалд хоёр үндсэн хэв
загвар хэрэглэгдэв:

\textbf{Синхрон харилцаа (Synchronous Communication):} HTTP/REST
протоколоор шууд дуудлага. \texttt{workflow\_service} нь статус
шилжилт хийхдээ \texttt{thesis\_service}-г шууд дуудаж дипломын
ажлын бичлэгийг шинэчлэнэ. Ийм шууд харилцааны тохиолдолд Circuit
Breaker загварыг (Nygard, 2007) хэрэглэж, upstream сервисийн
доройтол нь каскад байдлаар бусдад нэвтрэхээс хамгаалдаг.

\textbf{Асинхрон харилцаа (Asynchronous Communication):} Мэдэгдэл
(notification) явуулах, аудит бичлэг хийх гэх мэт цаг хугацааны
хамаарал бага үйлдлүүд нь Event-Driven загвараар хэрэгжих ёстой.
Ирээдүйн хөгжлийн шатанд Apache Kafka Message Broker нэвтрүүлснээр
\texttt{workflow\_service}-ийн статус шилжилтийн event нь
\texttt{notification\_service}-т асинхрон байдлаар хүрэх болно.

\begin{figure}[H]
    \centering
    \includegraphics[width=14cm]{fig_service_communication.png}
    \caption{Зураг 3.5. Микросервисүүдийн харилцааны диаграм. Хатуу шугам нь синхрон REST дуудлага, тасархай шугам нь асинхрон event-г тэмдэглэнэ. Эх сурвалж: зохиогчийн судалгаа, 2026.}
    \label{fig:service_communication}
\end{figure}

% ─────────────────────────────────────────────────────────────────────
\section{Өгөгдлийн Давхаргын Загварчлал}
% ─────────────────────────────────────────────────────────────────────

\subsection{Database Per Service загварын Сонголт}

Микросервисийн архитектурт өгөгдлийн давхаргыг хэрхэн хуваах вэ
гэсэн асуудал нь хамгийн нарийн төвөгтэй архитектурын шийдвэрүүдийн
нэг юм. Хоёр үндсэн сонголт байдаг:

\begin{enumerate}
  \item \textbf{Хуваалцсан өгөгдлийн сан (Shared Database):} Бүх
    сервис нэг өгөгдлийн санг хуваалцна. Хэрэгжүүлэлт энгийн боловч
    сервисүүдийн хооронд нягт холбоос (tight coupling) үүсдэг~---
    нэг сервисийн schema өөрчлөлт нь бусдад нөлөөлнө.
  \item \textbf{Сервис тус бүрийн өгөгдлийн сан (Database Per Service):}
    Сервис бүр өөрийн тусдаа schema эзэмшинэ. Сервис бүр бие даан
    deploy хийгдэх, технологийн сонголтоо чөлөөтэй хийх боломжтой
    боловч Cross-service join, distributed transaction нарийн болно.
\end{enumerate}

Энэхүү судалгаанд прагматик шийдлийг сонгосон: бүх сервис нэг
PostgreSQL instance-д ажилладаг боловч \textbf{schema-level тусгаарлалт}
хийгдсэн~--- сервис бүр өөрийн тусдаа schema-д бичиж, cross-schema
foreign key ашиглахгүй. Энэ нь ``Database Per Service'' загварын бие
даасан байдлын ашиг тусыг хэсэгчлэн хадгалахын зэрэгцээ нэг instance-ийн
хялбар засвар үйлчилгээний давуу талыг авч үлдэнэ.

\begin{figure}[H]
    \centering
    \includegraphics[width=14cm]{fig_database_schema_isolation.png}
    \caption{Зураг 3.6. Schema-level тусгаарлалтын загвар. Нэг PostgreSQL instance дотор сервис бүр өөрийн schema эзэмших бөгөөд хоорондын хамаарал нь application давхаргад шийдэгдэнэ. Эх сурвалж: зохиогчийн судалгаа, 2026.}
    \label{fig:db_schema_isolation}
\end{figure}

\subsection{Домейн Загварын Тодорхойлолт}

Системийн гол домейн объектуудын хоорондын харилцааг Entity-Relationship
(ER) загварт тодорхойлов. Хэдийгээр Cross-service join ашиглахгүй ч,
бизнесийн логикийн хэрэгцээгээр тодорхой домейн хооронд нягт харилцаа
оршино. Энэ харилцааг application давхаргад API дуудлагаар шийдвэрлэнэ.

Гол домейн объектууд болон тэдгээрийн атрибутуудыг дараах байдлаар
тодорхойлов:

\begin{itemize}
  \item \textbf{User} (user\_service): \texttt{id, sisiId, firstName,
    lastName, email, role, departmentId}
  \item \textbf{Department} (user\_service): \texttt{id, name, code,
    facultyId}
  \item \textbf{Topic} (topic\_service): \texttt{id, title, description,
    teacherId, maxStudents, status, academicYear}
  \item \textbf{Thesis} (thesis\_service): \texttt{id, title, studentId,
    teacherId, topicId, fileUrl, version, submittedAt}
  \item \textbf{WorkflowState} (workflow\_service): \texttt{id, thesisId,
    fromStatus, toStatus, actorId, comment, transitionedAt}
  \item \textbf{Committee} (committee\_service): \texttt{id, name,
    defenseDate, venueId, chairId}
  \item \textbf{CommitteeMember} (committee\_service): \texttt{id,
    committeeId, userId, role}
  \item \textbf{Evaluation} (evaluation\_service): \texttt{id, thesisId,
    committeeId, evaluatorId, score, criteria, submittedAt}
\end{itemize}

\begin{figure}[H]
    \centering
    \includegraphics[width=14cm]{fig_domain_er_diagram.png}
    \caption{Зураг 3.7. Системийн домейн загварын ER диаграм. Хатуу шугам нь нэг schema дотрох foreign key, тасархай шугам нь application давхаргын logical харилцааг заана. Эх сурвалж: зохиогчийн судалгаа, 2026.}
    \label{fig:domain_er}
\end{figure}

\subsection{Мэдээллийн Нийцэл ба Eventual Consistency}

Cross-service мэдээллийн нийцэлд ``Eventual Consistency'' (Аажмын
нийцэл) зарчмыг хэрэглэнэ. Distributed Transaction (2-Phase Commit)
нь микросервисийн орчинд хэдийгээр боломжтой боловч хүнд хэцүү
тогтолцоо шаарддаг тул Saga Pattern (Garcia-Molina \& Salem, 1987)
болон Outbox Pattern-ийг илүүд үзнэ.

Жишээ нь: оюутан дипломын ажлаа \texttt{SUBMITTED} төлөвт шилжүүлэхэд
\texttt{workflow\_service} нь шилжилтийн бичлэгийг хадгалж, дараа нь
\texttt{thesis\_service}-г шинэчлэх хүсэлт илгээнэ. Хэрэв
\texttt{thesis\_service} хариу өгөхгүй бол \texttt{workflow\_service}
нь retry логикоор дахин оролдох бөгөөд эцэст нь нийцэл тогтоно~---
энэ нь ``Eventual'' буюу аажмын нийцэл юм.

% ─────────────────────────────────────────────────────────────────────
\section{Нэгдсэн Архитектурын Загварын Нэгтгэл}
% ─────────────────────────────────────────────────────────────────────

Дээр тодорхойлсон гурван давхаргын загварчлал нэгдсэн байдлаар ажиллахад
дараах нэгдсэн зарчмуудыг дагана:

\begin{enumerate}
  \item \textbf{Нэг хариуцлага:} MFE модуль болон backend сервис тус
    бүр нэг бизнесийн чадварыг хариуцна~--- хариуцлагыг ``кэйс''
    дагуу биш, ``домейн'' дагуу хуваана.
  \item \textbf{API гэрээгээр харилцах:} Сервис болон MFE хоорондын
    харилцааны цорын ганц хэлбэр нь баримтжуулсан API гэрээ (OpenAPI
    Specification) юм~--- шууд function call, shared memory гэж байхгүй.
  \item \textbf{Бие даасан байршуулалт:} Сервис болон MFE тус бүр
    бусдыг зогсоолгүй шинэчлэгдэх чадвартай байна.
  \item \textbf{Барилгын бие даасан байдал:} Хөгжүүлэгч баг тус бүр
    өөрийн домейнийг бие даан эзэмшиж, технологийн сонголт дотор
    ухаалаг шийдэл гаргах эрхтэй.
\end{enumerate}

\begin{figure}[H]
    \centering
    \includegraphics[width=14cm]{fig_full_architecture_overview.png}
    \caption{Зураг 3.8. Системийн бүрэн архитектурын нэгдсэн бүдүүвч. Гурван давхарга (Frontend MFE, API Gateway, Backend Сервис) ба өгөгдлийн давхарга нэгдсэн байдлаар харагдана. Эх сурвалж: зохиогчийн судалгаа, 2026.}
    \label{fig:full_architecture}
\end{figure}

3-р бүлгийн загварчлалын шийдвэрүүд нь 4-р бүлгийн практик хэрэгжүүлэлтийн
суурь болно. Дараах бүлэгт ``юуг яагаад ийм загварчилсан'' гэдгийг
``хэрхэн бодитоор кодлон хэрэгжүүлсэн'' гэдэгтэй холбон тайлбарлана.
 гэж үндсэн файлд оруулна
% Шаардлагатай package-ууд: booktabs, float, graphicx, tabularx,
%   listings, amsmath, tikz (диаграмд), hyperref
% =====================================================================

\chapter{Системийн Архитектур ба Зохиомж}

Системийн архитектурын загварчлал нь зөвхөн технологийн сонголтын асуудал
биш~--- энэ нь байгууллагын хэрэгцээ, бизнесийн хил хязгаар, команд
бүтэц, ирээдүйн тэлэлтийн зорилтыг нийлүүлж, тогтвортой үндсэн шийдвэрт
хувиргах зохион байгуулалтын үйл явц юм. Энэхүү бүлэгт ``юуг яагаад
ингэж бүтэцлэх болсон'' гэсэн асуултад хариулж, системийн гурван үндсэн
давхарга~--- Frontend Microfrontend топологи, Backend Микросервисийн
топологи, Өгөгдлийн давхарга~--- тус бүрийн загварчлалын зарчим, хил
хязгаар, нэгдсэн дизайн шийдвэрийг дэлгэрэнгүй авч үзнэ.

% ─────────────────────────────────────────────────────────────────────
\section{Архитектурын Загварчлалын Үндэслэл}
% ─────────────────────────────────────────────────────────────────────

\subsection{Асуудлын Тодорхойлолт ба Архитектурын Шаардлага}

Монгол Улсын Их Сургуулийн дипломын ажлын удирдлагын одоогийн тогтолцоо
нь Jakarta Server Faces~2.2 (JSF) дээр суурилсан монолит вэб
программ хангамж бөгөөд бизнесийн логик, хуудасны загвар, мэдээллийн
давхарга нь нэгдсэн байдлаар хэрэгжсэн. Энэхүү монолит бүтэц нь системийн
хөгжлийн өмнө дараах тулгамдсан асуудлуудыг үүсгэж байв:

\begin{enumerate}
  \item \textbf{Масштабдах чадварын хязгаарлалт:} JVM heap дээрх
    HttpSession суурилсан тогтолцоо нь зэрэгцэх хэрэглэгчийн тоо
    нэмэгдэх тусам санах ойн хэрэглээ шугаман өсдөг тул өндөр
    ачааллын орчинд нарийн тохируулга шаарддаг.
  \item \textbf{Хөгжүүлэлтийн бие даасан байдлын дутагдал:} Нэг
    хуудасны өөрчлөлт нь бүтэн WAR байршуулалтыг шаарддаг тул
    frontend болон backend хөгжүүлэлт нэгэн зэрэг хийгдэх боломжгүй.
  \item \textbf{Орчин үеийн хэрэглэгчийн интерфэйс дутагдалтай байдал:}
    JSF-ийн request-response мөчлөг нь Single Page Application (SPA)-ийн
    responsive, real-time шинж чанарыг хязгаарладаг.
  \item \textbf{Технологийн хуучин өрийн хуримтлал:} JSF Facelets
    template-ийн нарийн шийдэл нь орчин үеийн TypeScript, React
    экосистемийн хөгжүүлэгчдэд суралцахад хэцүү тул элсүүлэлтийн
    зардал өндөр байдаг.
\end{enumerate}

Энэ асуудлуудад хариулах архитектурын үндсэн шаардлагуудыг Quality
Attribute Workshop (QAW, Barbacci et al., 2003) аргачлалаар тодорхойлж,
дараах таван тэргүүлэх шинж чанарыг сонгов: гүйцэтгэл, засвар
үйлчилгээний чадвар, бие даасан байршуулалтын чадвар, аюулгүй байдал,
шилжилтийн эрсдэлийн хамгаалалт.

\subsection{Архитектурын Стратегийн Сонголт: Аажмын Шилжилтийн Хандлага}

Одоо байгаа JSF системийг шилжүүлэхэд гурван стратегийн сонголт боломжтой:

\begin{enumerate}
  \item \textbf{Бүрэн Дахин Бичих (Big Bang Rewrite):} Одоогийн системийг
    бүрэн орхиж, шинийг эхнээс бичих. Онолын хувьд хамгийн цэвэр боловч
    Fowler (2004)-ийн тодорхойлсноор ``хамгийн эрсдэлтэй нэг юм бол
    бүхлийг нэгэн зэрэг дахин бичих'' учраас бизнесийн тасалдлыг дагуулна.
  \item \textbf{Хэв загварын Шинэчлэл (In-place Modernization):} JSF-ийг
    хадгалан, дотор нь орчин үеийн JavaScript фреймворкийг шингээх.
    Энэ нь JSF-ийн request-response мөчлөгийн хязгаарлалтыг арилгахгүй.
  \item \textbf{Strangler Fig загвар суурилсан Аажмын Шилжилт:} JSF-ийн
    монолитийг ``хост'' хэлбэрт хадгалж, шинэ React MFE модулиудыг
    аажмаар хоёр системийг зэрэгцүүлэн ажиллуулснаар нэгтгэх.
    Fowler (2004) болон Newman (2019)-ийн тодорхойлосон энэ загвар нь
    бизнесийн тасралтгүй байдлыг хамгийн сайн хангадаг.
\end{enumerate}

Энэхүү судалгаа нь гуравдахь сонголт~--- Strangler Fig загвар~--- г
сонгосон бөгөөд үндэслэл нь дараах байна: оюутны бүртгэл, дипломын
хамгаалалтын хуваарь, комиссын удирдлага зэрэг бизнесийн нарийн дараалалт
логикийг системийн ажиллагааг зогсоолгүйгээр шилжүүлэх шаардлага байсан.

\begin{table}[H]
    \centering
    \caption{Архитектурын стратегийн харьцуулалт}
    \label{tab:arch_strategy}
    \renewcommand{\arraystretch}{0.7}
    \begin{tabular}{lccc}
        \toprule
        \textbf{Шалгуур} & \textbf{Big Bang} & \textbf{In-place} & \textbf{Strangler Fig} \\
        \midrule
        Бизнесийн эрсдэл         & Өндөр  & Дунд   & Бага   \\
        Шилжилтийн хугацаа       & Богино & Урт    & Аажмын \\
        Технологийн цэвэр байдал & Өндөр  & Бага   & Дунд   \\
        Командын бие даасан байдал & Бага  & Бага   & Өндөр  \\
        Буцаалтын боломж         & Байхгүй & Хэцүү & Бүрэн  \\
        \bottomrule
    \end{tabular}
\end{table}

% ─────────────────────────────────────────────────────────────────────
\section{Frontend Microfrontend Архитектурын Загварчлал}
% ─────────────────────────────────────────────────────────────────────

\subsection{Гибрид Хостын Загварын Онолын Үндэслэл}

Microfrontend (MFE) архитектур нь backend микросервисийн зарчмуудыг
frontend давхаргад хэрэглэх ойлголт бөгөөд Cam Jackson (2019)-ийн
тодорхойлсноор ``Нэг хуудасны апп-ийг бие даасан, тусдаа байршуулагдах
жижиг фронтэнд апп-уудын нэгтгэл'' гэж тодорхойлогдоно.

Энэхүү системд JSF Tomcat нь нэгдсэн \textbf{Host} (Хост) буюу ``хаалт''
үүрэг гүйцэтгэж, React дээр суурилсан бие даасан модулиуд нь
\textbf{Remote} (Алслагдсан) MFE хэлбэрт ажиллана. Runtime Integration
(Ажиллах цагийн нэгтгэл) стратегийг ашигласан нь нэгтгэлийн цагт биш,
харин браузерт ачааллагдах үед модулиуд динамик байдлаар холбогдоно
гэсэн утгатай бөгөөд энэ нь Дэлхийн тэргүүлэх MFE архитектурын загвар
юм (Mezzalira, 2021).

\begin{figure}[H]
    \centering
    \includegraphics[width=14cm]{fig_hybrid_host_architecture.png}
    \caption{Зураг 3.1. Гибрид хостын загварын ерөнхий бүдүүвч. JSF Tomcat хост нь \texttt{/microfrontends/} замаар React Remote MFE-ийг ачааллаж, браузерт Runtime Integration хэрэгжинэ. Эх сурвалж: зохиогчийн судалгаа, 2026.}
    \label{fig:hybrid_host_arch}
\end{figure}

\subsection{Strangler Fig загварын Хэрэглэлт}

Martin Fowler (2004)-ийн Strangler Fig загвар нь Австралийн чулуун
инжирний мод (Ficus macrophylla)-ний органик загвараас санаа авсан:
шинэ ургамал нь хуучин мод дээр ургаж, аажмаар хучиж, хуучин мод
өөрөө унасны дараа бие даан зогсдог. Программ хангамжийн хүрээнд энэ нь:

\begin{enumerate}
  \item \textbf{Identify:} JSF-ийн аль хэсгийг хамгийн түрүүнд
    шилжүүлэх вэ гэдгийг бизнесийн үнэ цэнэ болон техникийн
    хамаарлын дагуу тодорхойлно.
  \item \textbf{Strangle:} Тухайн хэсгийг React MFE хэлбэрт хэрэгжүүлж,
    JSF-ийн хуудасны дотор ``цэнхэр нүх'' (injection point) хэлбэрт
    оруулна.
  \item \textbf{Kill:} Бүх хэсэг MFE болгон шилжсэний дараа JSF-ийн
    харгалзах хуудсыг устган, MFE-г бие даан ажиллуулна.
\end{enumerate}

Дипломын ажлын системд энэ гурван үе шат нь дараах дарааллаар хэрэгжсэн:
эхлээд нэвтрэлтийн хэлбэр (\texttt{module-auth}), дараа нь оюутны
ажлын хэсэг (\texttt{module-student}), дараа нь багшийн удирдлагын
хэсэг (\texttt{module-teacher}), эцэст нь хорооны хяналт
(\texttt{module-committee}).

\begin{figure}[H]
    \centering
    \includegraphics[width=14cm]{fig_strangler_fig_phases.png}
    \caption{Зураг 3.2. Strangler Fig загварын гурван үе шатны диаграм. Цагаан хэсэг нь JSF монолит, цэнхэр хэсэг нь React MFE, сумнууд нь шилжилтийн чиглэлийг заана. Эх сурвалж: зохиогчийн судалгаа, 2026.}
    \label{fig:strangler_phases}
\end{figure}

\subsection{MFE Модулиудын Загварчлал ба Хариуцлагын Хуваарилалт}

Системийн MFE бүтэц нь ``Bounded Context'' (Хязгаарлагдсан контекст,
Evans, 2003) зарчмын дагуу хуваагдав. Bounded Context нь домейны
мэдлэгийн хил хязгаарыг тодорхойлж, тус тусын домейны хэл, загвар,
хариуцлагыг тодорхой болгоно. Тус системд зургаан бие даасан MFE
модуль загварчлагдав:

\subsubsection{host-jsf (JSF Хост)}

JSF Tomcat сервер нь системийн нэгдсэн хаалт буюу \textbf{Application
Shell} үүрэг гүйцэтгэнэ. Үндсэн зорилго нь URL маршрутлалт, нэвтрэлтийн
хяналт, нийтлэг хуудасны бүтэц (header, navigation, footer) хадгалах,
мөн Remote MFE-ийн \texttt{remoteEntry.js} файлыг динамикаар ачаалах
injection точкийг хангах явдал юм. JSF хостын ``нүх'' бүр нь
\texttt{<div id=\texttt{"}microfrontend-root\texttt{"}/>} хэлбэртэй
DOM цэгт React-ийн render tree суурилна.

\subsubsection{module-auth (Нэвтрэлтийн MFE)}

Нэвтрэх хэлбэр, бүртгэлийн хэлбэр, нууц үг сэргээх урсгал зэргийг
хариуцна. Энэ модуль нь JWT токен авч, \texttt{localStorage}-д хадгалах
хариуцлага хүлээна. Бусад MFE-тэй харьцуулахад \texttt{module-auth} нь
бусад модулиудаас \textbf{тусгаарлагдсан} (isolated) ажиллах ёстой~---
нэвтрэлт хийгдэхэд хостын JSF хуудас руу чиглүүлэлт хийж, улмаар
хэрэглэгчийн дүрд тохирсон MFE ачааллагдана.

\subsubsection{module-student (Оюутны MFE)}

Оюутны хэрэглэгчийн интерфэйсийн бүрэн хариуцлагыг хүлээнэ: дипломын
сэдэв хүсэх, сэдвийн жагсаалт харах, удирдагч багштайгаа харилцах,
явцын тайлан оруулах, эцсийн бүтээлийн файл хавсаргах. Энэ модуль нь
\texttt{thesis\_service} болон \texttt{workflow\_service}-тэй API-аар
холбогдох бөгөөд real-time мэдэгдэл авах үүднээс
\texttt{notification\_service}-тэй WebSocket холболт тогтооно.

\subsubsection{module-teacher (Багшийн MFE)}

Удирдагч багшийн ажлын хэсгийг бүхэлд нь удирдана: оюутны дипломын
ажлыг зөвшөөрөх эсвэл буцаах, явцын тайлантай ажиллах, хэлэлцүүлгийн
хуваарь тохируулах, нийт ажлынхаа дашбоард харах. Багшийн MFE нь
\texttt{user\_service}-аас оюутны мэдээлэл татаж, \texttt{topic\_service}-аас
нэр дэвшигчийн жагсаалт авч ажиллана.

\subsubsection{module-committee (Хорооны MFE)}

Хамгаалалтын комиссын гишүүдэд зориулсан тусгай интерфэйс. Хорооны
бүрэлдэхүүн харах, хэлэлцүүлгийн хуваарийг баталгаажуулах, дүгнэлт
оруулах, нэгдсэн дүн гаргах үүрэгтэй. Энэ модуль нь ролын хязгаарлалт
хамгийн нарийн~--- зөвхөн \texttt{ROLE\_COMMITTEE\_MEMBER} дүр бүхий
хэрэглэгчид нэвтрэх боломжтой.

\subsubsection{module-thesis (Дипломын Ажлын Нэгдсэн MFE)}

Дипломын ажлын бүрэн амьдралын мөчлөгийн хяналт. Оюутан, багш,
хоёуланд нь харагдах нэгдсэн хэлбэрт харагддаг. Явцын Timeline,
баримт бичгийн хувилбарын түүх, тайлбарын гинж (comment thread)
зэргийг агуулна.

\subsubsection{shared-ui (Хуваалцсан UI Сан)}

Бусад MFE-д дундын хэрэгцээтэй бүрэлдэхүүнүүдийг нэгтгэн хадгалах
модуль. \texttt{PageHeader}, \texttt{StatusBadge}, \texttt{DataTable},
\texttt{PortalTabs} зэрэг бүрэлдэхүүнүүд нь энэ модульд хадгалагдаж,
бусад MFE-д \texttt{shared singleton} хэлбэрт нэгтгэгдэнэ. Ant Design
компонент санг тус модуль дамжуулан \textbf{нэгдсэн singleton} байдлаар
хуваалцах нь нийт bundle хэмжээг оновчтой байлгана.

\begin{figure}[H]
    \centering
    \includegraphics[width=14cm]{fig_mfe_module_map.png}
    \caption{Зураг 3.3. MFE модулиудын хариуцлагын хуваарилалт ба хоорондын харилцааны диаграм. Сум бүр нь ``хэрэглэнэ'' (consumes) харилцааг илэрхийлнэ. Эх сурвалж: зохиогчийн судалгаа, 2026.}
    \label{fig:mfe_module_map}
\end{figure}

\subsection{MFE Нэгтгэлийн Стратеги: Runtime Integration}

MFE нэгтгэлийн арга нь гурван ангилалд хуваагдана: \textbf{Build-time
Integration} (нэгтгэлийн үед), \textbf{Server-side Integration} (сервер
талын) болон \textbf{Runtime Integration} (ажиллах үеийн). Тус
системд Runtime Integration сонгосны шалтгаан нь:

\begin{itemize}
  \item Build-time Integration нь бие даасан байршуулалтыг зөвшөөрдөггүй~---
    нэг модуль өөрчлөгдөх бүрт бүгдийг дахин build хийх шаардлагатай болно.
  \item Server-side Integration (SSI) нь JSF-ийн request-response
    мөчлөгтэй зөрчилдөж, template бүрдэлтэд нарийн тохиргоо шаардана.
  \item Runtime Integration нь \textbf{Vite Module Federation} хэрэгслийн
    ачаар браузерт \texttt{remoteEntry.js} динамикаар татаж, зөвхөн
    шаардлагатай модулийг ачаалах ``lazy loading'' зарчмаар ажиллах
    боломжийг олгодог.
\end{itemize}

Runtime Integration-ийн загварт \texttt{remoteEntry.js} нь Remote MFE-ийн
``агуулгын хүснэгт'' (manifest) үүрэг гүйцэтгэж, Host нь энэ файлыг
ачаалснаар тухайн MFE-ийн экспортлосон бүрэлдэхүүнүүдийн мэдрэх зам
(resolution path)-ийг олж авдаг.

% ─────────────────────────────────────────────────────────────────────
\section{Backend Микросервисийн Архитектурын Загварчлал}
% ─────────────────────────────────────────────────────────────────────

\subsection{Монолитоос Микросервис Руу Шилжих Зарчим}

Backend-ийн архитектурын загварчлал нь Sam Newman (2019)-ийн
``Building Microservices'' болон Eric Evans (2003)-ийн Domain-Driven
Design зарчмуудад тулгуурлан хийгдсэн. Монолит backend нь ``Сервис бүр
өөрийн бие даасан байршуулалтын нэгж, нэг бизнесийн чадварыг хариуцна''
гэсэн Mikrosервисийн үндсэн зарчмын дагуу хуваагдсан.

Хуваах шалгуур нь дараах байв: нэгдүгээрт,
Хил хязгаарлагдсан контекст (Bounded Context)~--- домейны мэдлэгийн
хил хязгаар; хоёрдугаарт, Бизнесийн ажлын хэмжээний ялгаа~--- зарим
сервис нь өндөр ачааллыг дааж чаддаг байх шаардлагатай бол зарим нь
тийм ч их хэрэгцээгүй; гуравдугаарт, Командын хариуцлага~--- тус
командын тоо болон ур чадварын хуваарилалттай уялдуулах.

\begin{figure}[H]
    \centering
    \includegraphics[width=14cm]{fig_backend_microservices_overview.png}
    \caption{Зураг 3.4. Backend микросервисийн нэгдсэн бүдүүвч. API Gateway нь нэгдсэн нэвтрэх цэг болж, JWT баталгаажуулалт хийн, тохирох upstream сервис рүү дамжуулна. Эх сурвалж: зохиогчийн судалгаа, 2026.}
    \label{fig:backend_overview}
\end{figure}

\subsection{Микросервисүүдийн Тодорхойлолт ба Хариуцлага}

\subsubsection{auth-service (Баталгаажуулалтын Сервис)}

Хэрэглэгчийн нэвтрэлт, JWT токен үүсгэх, нэвтрэлтийн оролдлого
хянах, нууц үг сэргээх дараалал нь энэ сервисийн цөм хариуцлага юм.
Auth-service нь \texttt{num\_auth-main} багцын дор байрладаг бөгөөд
дотроо гурван дэд сервис агуулна: \texttt{auth-service} (нэвтрэлтийн
цөм), \texttt{authorization-service} (эрхийн удирдлага),
\texttt{user\_detail\_register} (хэрэглэгчийн дэлгэрэнгүй бүртгэл).

JWT-ийн загвар нь Stateless буюу ``Claims суурилсан'' загвар бөгөөд
токен нь \texttt{\{sisiId, roles, iat, exp\}} мэдэгдлийг агуулна.
Токений хэрэглэнгийн хугацаа 24 цаг байхаар тохируулагдсан бөгөөд
Refresh Token механизмийг нэмэлтээр дэмжнэ.

\subsubsection{user\_service (Хэрэглэгчийн Сервис)}

Хэрэглэгчийн профайлын бүрэн мэдрэх хариуцлагыг хүлээнэ: оюутан,
багш, хорооны гишүүн, хэлтсийн мэдээлэл. Auth-service нь зөвхөн
нэвтрэлтийн мэдлэгтэй (\texttt{sisiId}, нууц үг) харьцуулбал
user\_service нь дэлгэрэнгүй профайлын \textbf{единственный эх сурвалж}
(Single Source of Truth) болно. Сервисийн endpoints нь
\texttt{/api/users/students}, \texttt{/api/users/teachers},
\texttt{/api/users/by-email} зэрэг болно.

\subsubsection{topic\_service (Сэдвийн Сервис)}

Дипломын ажлын сэдвийн санд хариуцлагатай: сэдэв зарлах, хүсэлт
хүлээн авах, оюутан-багш хосолт батлах. Энэ сервис нь
\texttt{thesis\_service}-тэй нягт холбоотой боловч тусгаарлагдсан
байдаг~--- сэдвийн амьдралын мөчлөг нь дипломын ажлын амьдралын
мөчлөгтэй ялгаатай.

\subsubsection{thesis\_service (Дипломын Ажлын Сервис)}

Дипломын ажлын үндсэн мета мэдээллийг хадгалах, хувилбар удирдах,
баримт бичгийн хавсаргалт зохицуулах. Дипломын ажлын нэр, хаяг, хар
тэмдэглэл, файлын байршлын мета дата энд хадгалагдана. Файлын биет
агуулга нь object storage (эсвэл локал файлын систем)-д хадгалагдаж,
thesis\_service нь зөвхөн pointer хадгалдаг.

\subsubsection{workflow\_service (Явцын Удирдлагын Сервис)}

Дипломын ажлын амьдралын мөчлөгийн статус шилжилтийг удирдана.
State Machine загварт суурилсан бөгөөд төлөвүүд нь:
\texttt{DRAFT → SUBMITTED → UNDER\_REVIEW → REVISION\_REQUESTED →
APPROVED → DEFENSE\_SCHEDULED → DEFENDED → ARCHIVED}.
Шилжилт бүр нь хугацааны тэмдэглэл, хариуцлагатай хэрэглэгч,
тайлбарын бичлэгийг агуулна.

\subsubsection{evaluation\_service (Үнэлгээний Сервис)}

Комиссын гишүүн бүрийн дүгнэлтийг хүлээн авч, нэгтгэлийн тооцоолол
хийн, эцсийн дүн гаргах. Дүн гаргах алгоритм нь тохируулгатай
(configurable) байхаар загварчлагдсан~--- жишилтийн жин, нийт онооны
хязгаар гэх мэт.

\subsubsection{committee\_service (Комиссын Удирдлагын Сервис)}

Хэлэлцүүлгийн комиссын бүрэлдэхүүн тогтоох, хяналтын хуваарь
гаргах, комиссын гишүүдэд мэдэгдэл илгээх. Энэ сервис нь
\texttt{user\_service}-аас гишүүдийн мэдээлэл авч,
\texttt{notification\_service}-тэй хамтарч мэдэгдэл явуулна.

\subsubsection{API Gateway}

Нэгдсэн нэвтрэх цэг (Single Entry Point) үүрэг гүйцэтгэж, гадаад
хүсэлт бүрийг хүлээн авч, JWT шалгаж, тохирох upstream сервис рүү
дамжуулна (proxy/reverse proxy). Rate limiting, request logging,
CORS хяналт нь gateway давхаргад хэрэгжинэ тул backend сервис тус
бүр энэ логикийг давтан хэрэгжүүлэх шаардлагагүй болно.

\begin{table}[H]
    \centering
    \caption{Backend микросервисүүдийн порт болон хариуцлагын хуваарь}
    \label{tab:service_ports}
    \renewcommand{\arraystretch}{0.7}
    \begin{tabular}{llll}
        \toprule
        \textbf{Сервис} & \textbf{Порт} & \textbf{Технологи} & \textbf{Гол хариуцлага} \\
        \midrule
        auth-service         & 8887 & Spring Boot  & JWT үүсгэх, нэвтрэлт \\
        user\_service        & 8086 & Spring WebFlux & Хэрэглэгчийн профайл \\
        topic\_service       & 8081 & Spring WebFlux & Сэдвийн удирдлага \\
        thesis\_service      & 8083 & Spring WebFlux & Дипломын ажлын мета \\
        workflow\_service    & 8084 & Spring WebFlux & Статус шилжилт \\
        evaluation\_service  & 8085 & Spring WebFlux & Дүгнэлт, үнэлгээ \\
        committee\_service   & 8082 & Spring WebFlux & Комиссын удирдлага \\
        Tomcat (JSF+MFE)     & 8080 & Apache Tomcat  & Host, static file serve \\
        \bottomrule
    \end{tabular}
\end{table}

\subsection{Сервисүүдийн Хоорондын Харилцааны Загвар}

Микросервисүүдийн хоорондын харилцааны загварчлалд хоёр үндсэн хэв
загвар хэрэглэгдэв:

\textbf{Синхрон харилцаа (Synchronous Communication):} HTTP/REST
протоколоор шууд дуудлага. \texttt{workflow\_service} нь статус
шилжилт хийхдээ \texttt{thesis\_service}-г шууд дуудаж дипломын
ажлын бичлэгийг шинэчлэнэ. Ийм шууд харилцааны тохиолдолд Circuit
Breaker загварыг (Nygard, 2007) хэрэглэж, upstream сервисийн
доройтол нь каскад байдлаар бусдад нэвтрэхээс хамгаалдаг.

\textbf{Асинхрон харилцаа (Asynchronous Communication):} Мэдэгдэл
(notification) явуулах, аудит бичлэг хийх гэх мэт цаг хугацааны
хамаарал бага үйлдлүүд нь Event-Driven загвараар хэрэгжих ёстой.
Ирээдүйн хөгжлийн шатанд Apache Kafka Message Broker нэвтрүүлснээр
\texttt{workflow\_service}-ийн статус шилжилтийн event нь
\texttt{notification\_service}-т асинхрон байдлаар хүрэх болно.

\begin{figure}[H]
    \centering
    \includegraphics[width=14cm]{fig_service_communication.png}
    \caption{Зураг 3.5. Микросервисүүдийн харилцааны диаграм. Хатуу шугам нь синхрон REST дуудлага, тасархай шугам нь асинхрон event-г тэмдэглэнэ. Эх сурвалж: зохиогчийн судалгаа, 2026.}
    \label{fig:service_communication}
\end{figure}

% ─────────────────────────────────────────────────────────────────────
\section{Өгөгдлийн Давхаргын Загварчлал}
% ─────────────────────────────────────────────────────────────────────

\subsection{Database Per Service загварын Сонголт}

Микросервисийн архитектурт өгөгдлийн давхаргыг хэрхэн хуваах вэ
гэсэн асуудал нь хамгийн нарийн төвөгтэй архитектурын шийдвэрүүдийн
нэг юм. Хоёр үндсэн сонголт байдаг:

\begin{enumerate}
  \item \textbf{Хуваалцсан өгөгдлийн сан (Shared Database):} Бүх
    сервис нэг өгөгдлийн санг хуваалцна. Хэрэгжүүлэлт энгийн боловч
    сервисүүдийн хооронд нягт холбоос (tight coupling) үүсдэг~---
    нэг сервисийн schema өөрчлөлт нь бусдад нөлөөлнө.
  \item \textbf{Сервис тус бүрийн өгөгдлийн сан (Database Per Service):}
    Сервис бүр өөрийн тусдаа schema эзэмшинэ. Сервис бүр бие даан
    deploy хийгдэх, технологийн сонголтоо чөлөөтэй хийх боломжтой
    боловч Cross-service join, distributed transaction нарийн болно.
\end{enumerate}

Энэхүү судалгаанд прагматик шийдлийг сонгосон: бүх сервис нэг
PostgreSQL instance-д ажилладаг боловч \textbf{schema-level тусгаарлалт}
хийгдсэн~--- сервис бүр өөрийн тусдаа schema-д бичиж, cross-schema
foreign key ашиглахгүй. Энэ нь ``Database Per Service'' загварын бие
даасан байдлын ашиг тусыг хэсэгчлэн хадгалахын зэрэгцээ нэг instance-ийн
хялбар засвар үйлчилгээний давуу талыг авч үлдэнэ.

\begin{figure}[H]
    \centering
    \includegraphics[width=14cm]{fig_database_schema_isolation.png}
    \caption{Зураг 3.6. Schema-level тусгаарлалтын загвар. Нэг PostgreSQL instance дотор сервис бүр өөрийн schema эзэмших бөгөөд хоорондын хамаарал нь application давхаргад шийдэгдэнэ. Эх сурвалж: зохиогчийн судалгаа, 2026.}
    \label{fig:db_schema_isolation}
\end{figure}

\subsection{Домейн Загварын Тодорхойлолт}

Системийн гол домейн объектуудын хоорондын харилцааг Entity-Relationship
(ER) загварт тодорхойлов. Хэдийгээр Cross-service join ашиглахгүй ч,
бизнесийн логикийн хэрэгцээгээр тодорхой домейн хооронд нягт харилцаа
оршино. Энэ харилцааг application давхаргад API дуудлагаар шийдвэрлэнэ.

Гол домейн объектууд болон тэдгээрийн атрибутуудыг дараах байдлаар
тодорхойлов:

\begin{itemize}
  \item \textbf{User} (user\_service): \texttt{id, sisiId, firstName,
    lastName, email, role, departmentId}
  \item \textbf{Department} (user\_service): \texttt{id, name, code,
    facultyId}
  \item \textbf{Topic} (topic\_service): \texttt{id, title, description,
    teacherId, maxStudents, status, academicYear}
  \item \textbf{Thesis} (thesis\_service): \texttt{id, title, studentId,
    teacherId, topicId, fileUrl, version, submittedAt}
  \item \textbf{WorkflowState} (workflow\_service): \texttt{id, thesisId,
    fromStatus, toStatus, actorId, comment, transitionedAt}
  \item \textbf{Committee} (committee\_service): \texttt{id, name,
    defenseDate, venueId, chairId}
  \item \textbf{CommitteeMember} (committee\_service): \texttt{id,
    committeeId, userId, role}
  \item \textbf{Evaluation} (evaluation\_service): \texttt{id, thesisId,
    committeeId, evaluatorId, score, criteria, submittedAt}
\end{itemize}

\begin{figure}[H]
    \centering
    \includegraphics[width=14cm]{fig_domain_er_diagram.png}
    \caption{Зураг 3.7. Системийн домейн загварын ER диаграм. Хатуу шугам нь нэг schema дотрох foreign key, тасархай шугам нь application давхаргын logical харилцааг заана. Эх сурвалж: зохиогчийн судалгаа, 2026.}
    \label{fig:domain_er}
\end{figure}

\subsection{Мэдээллийн Нийцэл ба Eventual Consistency}

Cross-service мэдээллийн нийцэлд ``Eventual Consistency'' (Аажмын
нийцэл) зарчмыг хэрэглэнэ. Distributed Transaction (2-Phase Commit)
нь микросервисийн орчинд хэдийгээр боломжтой боловч хүнд хэцүү
тогтолцоо шаарддаг тул Saga Pattern (Garcia-Molina \& Salem, 1987)
болон Outbox Pattern-ийг илүүд үзнэ.

Жишээ нь: оюутан дипломын ажлаа \texttt{SUBMITTED} төлөвт шилжүүлэхэд
\texttt{workflow\_service} нь шилжилтийн бичлэгийг хадгалж, дараа нь
\texttt{thesis\_service}-г шинэчлэх хүсэлт илгээнэ. Хэрэв
\texttt{thesis\_service} хариу өгөхгүй бол \texttt{workflow\_service}
нь retry логикоор дахин оролдох бөгөөд эцэст нь нийцэл тогтоно~---
энэ нь ``Eventual'' буюу аажмын нийцэл юм.

% ─────────────────────────────────────────────────────────────────────
\section{Нэгдсэн Архитектурын Загварын Нэгтгэл}
% ─────────────────────────────────────────────────────────────────────

Дээр тодорхойлсон гурван давхаргын загварчлал нэгдсэн байдлаар ажиллахад
дараах нэгдсэн зарчмуудыг дагана:

\begin{enumerate}
  \item \textbf{Нэг хариуцлага:} MFE модуль болон backend сервис тус
    бүр нэг бизнесийн чадварыг хариуцна~--- хариуцлагыг ``кэйс''
    дагуу биш, ``домейн'' дагуу хуваана.
  \item \textbf{API гэрээгээр харилцах:} Сервис болон MFE хоорондын
    харилцааны цорын ганц хэлбэр нь баримтжуулсан API гэрээ (OpenAPI
    Specification) юм~--- шууд function call, shared memory гэж байхгүй.
  \item \textbf{Бие даасан байршуулалт:} Сервис болон MFE тус бүр
    бусдыг зогсоолгүй шинэчлэгдэх чадвартай байна.
  \item \textbf{Барилгын бие даасан байдал:} Хөгжүүлэгч баг тус бүр
    өөрийн домейнийг бие даан эзэмшиж, технологийн сонголт дотор
    ухаалаг шийдэл гаргах эрхтэй.
\end{enumerate}

\begin{figure}[H]
    \centering
    \includegraphics[width=14cm]{fig_full_architecture_overview.png}
    \caption{Зураг 3.8. Системийн бүрэн архитектурын нэгдсэн бүдүүвч. Гурван давхарга (Frontend MFE, API Gateway, Backend Сервис) ба өгөгдлийн давхарга нэгдсэн байдлаар харагдана. Эх сурвалж: зохиогчийн судалгаа, 2026.}
    \label{fig:full_architecture}
\end{figure}

3-р бүлгийн загварчлалын шийдвэрүүд нь 4-р бүлгийн практик хэрэгжүүлэлтийн
суурь болно. Дараах бүлэгт ``юуг яагаад ийм загварчилсан'' гэдгийг
``хэрхэн бодитоор кодлон хэрэгжүүлсэн'' гэдэгтэй холбон тайлбарлана.
 гэж үндсэн файлд оруулна
% Шаардлагатай package-ууд: booktabs, float, graphicx, tabularx,
%   listings, amsmath, tikz (диаграмд), hyperref
% =====================================================================

\chapter{Системийн Архитектур ба Зохиомж}

Системийн архитектурын загварчлал нь зөвхөн технологийн сонголтын асуудал
биш~--- энэ нь байгууллагын хэрэгцээ, бизнесийн хил хязгаар, команд
бүтэц, ирээдүйн тэлэлтийн зорилтыг нийлүүлж, тогтвортой үндсэн шийдвэрт
хувиргах зохион байгуулалтын үйл явц юм. Энэхүү бүлэгт ``юуг яагаад
ингэж бүтэцлэх болсон'' гэсэн асуултад хариулж, системийн гурван үндсэн
давхарга~--- Frontend Microfrontend топологи, Backend Микросервисийн
топологи, Өгөгдлийн давхарга~--- тус бүрийн загварчлалын зарчим, хил
хязгаар, нэгдсэн дизайн шийдвэрийг дэлгэрэнгүй авч үзнэ.

% ─────────────────────────────────────────────────────────────────────
\section{Архитектурын Загварчлалын Үндэслэл}
% ─────────────────────────────────────────────────────────────────────

\subsection{Асуудлын Тодорхойлолт ба Архитектурын Шаардлага}

Монгол Улсын Их Сургуулийн дипломын ажлын удирдлагын одоогийн тогтолцоо
нь Jakarta Server Faces~2.2 (JSF) дээр суурилсан монолит вэб
программ хангамж бөгөөд бизнесийн логик, хуудасны загвар, мэдээллийн
давхарга нь нэгдсэн байдлаар хэрэгжсэн. Энэхүү монолит бүтэц нь системийн
хөгжлийн өмнө дараах тулгамдсан асуудлуудыг үүсгэж байв:

\begin{enumerate}
  \item \textbf{Масштабдах чадварын хязгаарлалт:} JVM heap дээрх
    HttpSession суурилсан тогтолцоо нь зэрэгцэх хэрэглэгчийн тоо
    нэмэгдэх тусам санах ойн хэрэглээ шугаман өсдөг тул өндөр
    ачааллын орчинд нарийн тохируулга шаарддаг.
  \item \textbf{Хөгжүүлэлтийн бие даасан байдлын дутагдал:} Нэг
    хуудасны өөрчлөлт нь бүтэн WAR байршуулалтыг шаарддаг тул
    frontend болон backend хөгжүүлэлт нэгэн зэрэг хийгдэх боломжгүй.
  \item \textbf{Орчин үеийн хэрэглэгчийн интерфэйс дутагдалтай байдал:}
    JSF-ийн request-response мөчлөг нь Single Page Application (SPA)-ийн
    responsive, real-time шинж чанарыг хязгаарладаг.
  \item \textbf{Технологийн хуучин өрийн хуримтлал:} JSF Facelets
    template-ийн нарийн шийдэл нь орчин үеийн TypeScript, React
    экосистемийн хөгжүүлэгчдэд суралцахад хэцүү тул элсүүлэлтийн
    зардал өндөр байдаг.
\end{enumerate}

Энэ асуудлуудад хариулах архитектурын үндсэн шаардлагуудыг Quality
Attribute Workshop (QAW, Barbacci et al., 2003) аргачлалаар тодорхойлж,
дараах таван тэргүүлэх шинж чанарыг сонгов: гүйцэтгэл, засвар
үйлчилгээний чадвар, бие даасан байршуулалтын чадвар, аюулгүй байдал,
шилжилтийн эрсдэлийн хамгаалалт.

\subsection{Архитектурын Стратегийн Сонголт: Аажмын Шилжилтийн Хандлага}

Одоо байгаа JSF системийг шилжүүлэхэд гурван стратегийн сонголт боломжтой:

\begin{enumerate}
  \item \textbf{Бүрэн Дахин Бичих (Big Bang Rewrite):} Одоогийн системийг
    бүрэн орхиж, шинийг эхнээс бичих. Онолын хувьд хамгийн цэвэр боловч
    Fowler (2004)-ийн тодорхойлсноор ``хамгийн эрсдэлтэй нэг юм бол
    бүхлийг нэгэн зэрэг дахин бичих'' учраас бизнесийн тасалдлыг дагуулна.
  \item \textbf{Хэв загварын Шинэчлэл (In-place Modernization):} JSF-ийг
    хадгалан, дотор нь орчин үеийн JavaScript фреймворкийг шингээх.
    Энэ нь JSF-ийн request-response мөчлөгийн хязгаарлалтыг арилгахгүй.
  \item \textbf{Strangler Fig загвар суурилсан Аажмын Шилжилт:} JSF-ийн
    монолитийг ``хост'' хэлбэрт хадгалж, шинэ React MFE модулиудыг
    аажмаар хоёр системийг зэрэгцүүлэн ажиллуулснаар нэгтгэх.
    Fowler (2004) болон Newman (2019)-ийн тодорхойлосон энэ загвар нь
    бизнесийн тасралтгүй байдлыг хамгийн сайн хангадаг.
\end{enumerate}

Энэхүү судалгаа нь гуравдахь сонголт~--- Strangler Fig загвар~--- г
сонгосон бөгөөд үндэслэл нь дараах байна: оюутны бүртгэл, дипломын
хамгаалалтын хуваарь, комиссын удирдлага зэрэг бизнесийн нарийн дараалалт
логикийг системийн ажиллагааг зогсоолгүйгээр шилжүүлэх шаардлага байсан.

\begin{table}[H]
    \centering
    \caption{Архитектурын стратегийн харьцуулалт}
    \label{tab:arch_strategy}
    \renewcommand{\arraystretch}{0.7}
    \begin{tabular}{lccc}
        \toprule
        \textbf{Шалгуур} & \textbf{Big Bang} & \textbf{In-place} & \textbf{Strangler Fig} \\
        \midrule
        Бизнесийн эрсдэл         & Өндөр  & Дунд   & Бага   \\
        Шилжилтийн хугацаа       & Богино & Урт    & Аажмын \\
        Технологийн цэвэр байдал & Өндөр  & Бага   & Дунд   \\
        Командын бие даасан байдал & Бага  & Бага   & Өндөр  \\
        Буцаалтын боломж         & Байхгүй & Хэцүү & Бүрэн  \\
        \bottomrule
    \end{tabular}
\end{table}

% ─────────────────────────────────────────────────────────────────────
\section{Frontend Microfrontend Архитектурын Загварчлал}
% ─────────────────────────────────────────────────────────────────────

\subsection{Гибрид Хостын Загварын Онолын Үндэслэл}

Microfrontend (MFE) архитектур нь backend микросервисийн зарчмуудыг
frontend давхаргад хэрэглэх ойлголт бөгөөд Cam Jackson (2019)-ийн
тодорхойлсноор ``Нэг хуудасны апп-ийг бие даасан, тусдаа байршуулагдах
жижиг фронтэнд апп-уудын нэгтгэл'' гэж тодорхойлогдоно.

Энэхүү системд JSF Tomcat нь нэгдсэн \textbf{Host} (Хост) буюу ``хаалт''
үүрэг гүйцэтгэж, React дээр суурилсан бие даасан модулиуд нь
\textbf{Remote} (Алслагдсан) MFE хэлбэрт ажиллана. Runtime Integration
(Ажиллах цагийн нэгтгэл) стратегийг ашигласан нь нэгтгэлийн цагт биш,
харин браузерт ачааллагдах үед модулиуд динамик байдлаар холбогдоно
гэсэн утгатай бөгөөд энэ нь Дэлхийн тэргүүлэх MFE архитектурын загвар
юм (Mezzalira, 2021).

\begin{figure}[H]
    \centering
    \includegraphics[width=14cm]{fig_hybrid_host_architecture.png}
    \caption{Зураг 3.1. Гибрид хостын загварын ерөнхий бүдүүвч. JSF Tomcat хост нь \texttt{/microfrontends/} замаар React Remote MFE-ийг ачааллаж, браузерт Runtime Integration хэрэгжинэ. Эх сурвалж: зохиогчийн судалгаа, 2026.}
    \label{fig:hybrid_host_arch}
\end{figure}

\subsection{Strangler Fig загварын Хэрэглэлт}

Martin Fowler (2004)-ийн Strangler Fig загвар нь Австралийн чулуун
инжирний мод (Ficus macrophylla)-ний органик загвараас санаа авсан:
шинэ ургамал нь хуучин мод дээр ургаж, аажмаар хучиж, хуучин мод
өөрөө унасны дараа бие даан зогсдог. Программ хангамжийн хүрээнд энэ нь:

\begin{enumerate}
  \item \textbf{Identify:} JSF-ийн аль хэсгийг хамгийн түрүүнд
    шилжүүлэх вэ гэдгийг бизнесийн үнэ цэнэ болон техникийн
    хамаарлын дагуу тодорхойлно.
  \item \textbf{Strangle:} Тухайн хэсгийг React MFE хэлбэрт хэрэгжүүлж,
    JSF-ийн хуудасны дотор ``цэнхэр нүх'' (injection point) хэлбэрт
    оруулна.
  \item \textbf{Kill:} Бүх хэсэг MFE болгон шилжсэний дараа JSF-ийн
    харгалзах хуудсыг устган, MFE-г бие даан ажиллуулна.
\end{enumerate}

Дипломын ажлын системд энэ гурван үе шат нь дараах дарааллаар хэрэгжсэн:
эхлээд нэвтрэлтийн хэлбэр (\texttt{module-auth}), дараа нь оюутны
ажлын хэсэг (\texttt{module-student}), дараа нь багшийн удирдлагын
хэсэг (\texttt{module-teacher}), эцэст нь хорооны хяналт
(\texttt{module-committee}).

\begin{figure}[H]
    \centering
    \includegraphics[width=14cm]{fig_strangler_fig_phases.png}
    \caption{Зураг 3.2. Strangler Fig загварын гурван үе шатны диаграм. Цагаан хэсэг нь JSF монолит, цэнхэр хэсэг нь React MFE, сумнууд нь шилжилтийн чиглэлийг заана. Эх сурвалж: зохиогчийн судалгаа, 2026.}
    \label{fig:strangler_phases}
\end{figure}

\subsection{MFE Модулиудын Загварчлал ба Хариуцлагын Хуваарилалт}

Системийн MFE бүтэц нь ``Bounded Context'' (Хязгаарлагдсан контекст,
Evans, 2003) зарчмын дагуу хуваагдав. Bounded Context нь домейны
мэдлэгийн хил хязгаарыг тодорхойлж, тус тусын домейны хэл, загвар,
хариуцлагыг тодорхой болгоно. Тус системд зургаан бие даасан MFE
модуль загварчлагдав:

\subsubsection{host-jsf (JSF Хост)}

JSF Tomcat сервер нь системийн нэгдсэн хаалт буюу \textbf{Application
Shell} үүрэг гүйцэтгэнэ. Үндсэн зорилго нь URL маршрутлалт, нэвтрэлтийн
хяналт, нийтлэг хуудасны бүтэц (header, navigation, footer) хадгалах,
мөн Remote MFE-ийн \texttt{remoteEntry.js} файлыг динамикаар ачаалах
injection точкийг хангах явдал юм. JSF хостын ``нүх'' бүр нь
\texttt{<div id=\texttt{"}microfrontend-root\texttt{"}/>} хэлбэртэй
DOM цэгт React-ийн render tree суурилна.

\subsubsection{module-auth (Нэвтрэлтийн MFE)}

Нэвтрэх хэлбэр, бүртгэлийн хэлбэр, нууц үг сэргээх урсгал зэргийг
хариуцна. Энэ модуль нь JWT токен авч, \texttt{localStorage}-д хадгалах
хариуцлага хүлээна. Бусад MFE-тэй харьцуулахад \texttt{module-auth} нь
бусад модулиудаас \textbf{тусгаарлагдсан} (isolated) ажиллах ёстой~---
нэвтрэлт хийгдэхэд хостын JSF хуудас руу чиглүүлэлт хийж, улмаар
хэрэглэгчийн дүрд тохирсон MFE ачааллагдана.

\subsubsection{module-student (Оюутны MFE)}

Оюутны хэрэглэгчийн интерфэйсийн бүрэн хариуцлагыг хүлээнэ: дипломын
сэдэв хүсэх, сэдвийн жагсаалт харах, удирдагч багштайгаа харилцах,
явцын тайлан оруулах, эцсийн бүтээлийн файл хавсаргах. Энэ модуль нь
\texttt{thesis\_service} болон \texttt{workflow\_service}-тэй API-аар
холбогдох бөгөөд real-time мэдэгдэл авах үүднээс
\texttt{notification\_service}-тэй WebSocket холболт тогтооно.

\subsubsection{module-teacher (Багшийн MFE)}

Удирдагч багшийн ажлын хэсгийг бүхэлд нь удирдана: оюутны дипломын
ажлыг зөвшөөрөх эсвэл буцаах, явцын тайлантай ажиллах, хэлэлцүүлгийн
хуваарь тохируулах, нийт ажлынхаа дашбоард харах. Багшийн MFE нь
\texttt{user\_service}-аас оюутны мэдээлэл татаж, \texttt{topic\_service}-аас
нэр дэвшигчийн жагсаалт авч ажиллана.

\subsubsection{module-committee (Хорооны MFE)}

Хамгаалалтын комиссын гишүүдэд зориулсан тусгай интерфэйс. Хорооны
бүрэлдэхүүн харах, хэлэлцүүлгийн хуваарийг баталгаажуулах, дүгнэлт
оруулах, нэгдсэн дүн гаргах үүрэгтэй. Энэ модуль нь ролын хязгаарлалт
хамгийн нарийн~--- зөвхөн \texttt{ROLE\_COMMITTEE\_MEMBER} дүр бүхий
хэрэглэгчид нэвтрэх боломжтой.

\subsubsection{module-thesis (Дипломын Ажлын Нэгдсэн MFE)}

Дипломын ажлын бүрэн амьдралын мөчлөгийн хяналт. Оюутан, багш,
хоёуланд нь харагдах нэгдсэн хэлбэрт харагддаг. Явцын Timeline,
баримт бичгийн хувилбарын түүх, тайлбарын гинж (comment thread)
зэргийг агуулна.

\subsubsection{shared-ui (Хуваалцсан UI Сан)}

Бусад MFE-д дундын хэрэгцээтэй бүрэлдэхүүнүүдийг нэгтгэн хадгалах
модуль. \texttt{PageHeader}, \texttt{StatusBadge}, \texttt{DataTable},
\texttt{PortalTabs} зэрэг бүрэлдэхүүнүүд нь энэ модульд хадгалагдаж,
бусад MFE-д \texttt{shared singleton} хэлбэрт нэгтгэгдэнэ. Ant Design
компонент санг тус модуль дамжуулан \textbf{нэгдсэн singleton} байдлаар
хуваалцах нь нийт bundle хэмжээг оновчтой байлгана.

\begin{figure}[H]
    \centering
    \includegraphics[width=14cm]{fig_mfe_module_map.png}
    \caption{Зураг 3.3. MFE модулиудын хариуцлагын хуваарилалт ба хоорондын харилцааны диаграм. Сум бүр нь ``хэрэглэнэ'' (consumes) харилцааг илэрхийлнэ. Эх сурвалж: зохиогчийн судалгаа, 2026.}
    \label{fig:mfe_module_map}
\end{figure}

\subsection{MFE Нэгтгэлийн Стратеги: Runtime Integration}

MFE нэгтгэлийн арга нь гурван ангилалд хуваагдана: \textbf{Build-time
Integration} (нэгтгэлийн үед), \textbf{Server-side Integration} (сервер
талын) болон \textbf{Runtime Integration} (ажиллах үеийн). Тус
системд Runtime Integration сонгосны шалтгаан нь:

\begin{itemize}
  \item Build-time Integration нь бие даасан байршуулалтыг зөвшөөрдөггүй~---
    нэг модуль өөрчлөгдөх бүрт бүгдийг дахин build хийх шаардлагатай болно.
  \item Server-side Integration (SSI) нь JSF-ийн request-response
    мөчлөгтэй зөрчилдөж, template бүрдэлтэд нарийн тохиргоо шаардана.
  \item Runtime Integration нь \textbf{Vite Module Federation} хэрэгслийн
    ачаар браузерт \texttt{remoteEntry.js} динамикаар татаж, зөвхөн
    шаардлагатай модулийг ачаалах ``lazy loading'' зарчмаар ажиллах
    боломжийг олгодог.
\end{itemize}

Runtime Integration-ийн загварт \texttt{remoteEntry.js} нь Remote MFE-ийн
``агуулгын хүснэгт'' (manifest) үүрэг гүйцэтгэж, Host нь энэ файлыг
ачаалснаар тухайн MFE-ийн экспортлосон бүрэлдэхүүнүүдийн мэдрэх зам
(resolution path)-ийг олж авдаг.

% ─────────────────────────────────────────────────────────────────────
\section{Backend Микросервисийн Архитектурын Загварчлал}
% ─────────────────────────────────────────────────────────────────────

\subsection{Монолитоос Микросервис Руу Шилжих Зарчим}

Backend-ийн архитектурын загварчлал нь Sam Newman (2019)-ийн
``Building Microservices'' болон Eric Evans (2003)-ийн Domain-Driven
Design зарчмуудад тулгуурлан хийгдсэн. Монолит backend нь ``Сервис бүр
өөрийн бие даасан байршуулалтын нэгж, нэг бизнесийн чадварыг хариуцна''
гэсэн Mikrosервисийн үндсэн зарчмын дагуу хуваагдсан.

Хуваах шалгуур нь дараах байв: нэгдүгээрт,
Хил хязгаарлагдсан контекст (Bounded Context)~--- домейны мэдлэгийн
хил хязгаар; хоёрдугаарт, Бизнесийн ажлын хэмжээний ялгаа~--- зарим
сервис нь өндөр ачааллыг дааж чаддаг байх шаардлагатай бол зарим нь
тийм ч их хэрэгцээгүй; гуравдугаарт, Командын хариуцлага~--- тус
командын тоо болон ур чадварын хуваарилалттай уялдуулах.

\begin{figure}[H]
    \centering
    \includegraphics[width=14cm]{fig_backend_microservices_overview.png}
    \caption{Зураг 3.4. Backend микросервисийн нэгдсэн бүдүүвч. API Gateway нь нэгдсэн нэвтрэх цэг болж, JWT баталгаажуулалт хийн, тохирох upstream сервис рүү дамжуулна. Эх сурвалж: зохиогчийн судалгаа, 2026.}
    \label{fig:backend_overview}
\end{figure}

\subsection{Микросервисүүдийн Тодорхойлолт ба Хариуцлага}

\subsubsection{auth-service (Баталгаажуулалтын Сервис)}

Хэрэглэгчийн нэвтрэлт, JWT токен үүсгэх, нэвтрэлтийн оролдлого
хянах, нууц үг сэргээх дараалал нь энэ сервисийн цөм хариуцлага юм.
Auth-service нь \texttt{num\_auth-main} багцын дор байрладаг бөгөөд
дотроо гурван дэд сервис агуулна: \texttt{auth-service} (нэвтрэлтийн
цөм), \texttt{authorization-service} (эрхийн удирдлага),
\texttt{user\_detail\_register} (хэрэглэгчийн дэлгэрэнгүй бүртгэл).

JWT-ийн загвар нь Stateless буюу ``Claims суурилсан'' загвар бөгөөд
токен нь \texttt{\{sisiId, roles, iat, exp\}} мэдэгдлийг агуулна.
Токений хэрэглэнгийн хугацаа 24 цаг байхаар тохируулагдсан бөгөөд
Refresh Token механизмийг нэмэлтээр дэмжнэ.

\subsubsection{user\_service (Хэрэглэгчийн Сервис)}

Хэрэглэгчийн профайлын бүрэн мэдрэх хариуцлагыг хүлээнэ: оюутан,
багш, хорооны гишүүн, хэлтсийн мэдээлэл. Auth-service нь зөвхөн
нэвтрэлтийн мэдлэгтэй (\texttt{sisiId}, нууц үг) харьцуулбал
user\_service нь дэлгэрэнгүй профайлын \textbf{единственный эх сурвалж}
(Single Source of Truth) болно. Сервисийн endpoints нь
\texttt{/api/users/students}, \texttt{/api/users/teachers},
\texttt{/api/users/by-email} зэрэг болно.

\subsubsection{topic\_service (Сэдвийн Сервис)}

Дипломын ажлын сэдвийн санд хариуцлагатай: сэдэв зарлах, хүсэлт
хүлээн авах, оюутан-багш хосолт батлах. Энэ сервис нь
\texttt{thesis\_service}-тэй нягт холбоотой боловч тусгаарлагдсан
байдаг~--- сэдвийн амьдралын мөчлөг нь дипломын ажлын амьдралын
мөчлөгтэй ялгаатай.

\subsubsection{thesis\_service (Дипломын Ажлын Сервис)}

Дипломын ажлын үндсэн мета мэдээллийг хадгалах, хувилбар удирдах,
баримт бичгийн хавсаргалт зохицуулах. Дипломын ажлын нэр, хаяг, хар
тэмдэглэл, файлын байршлын мета дата энд хадгалагдана. Файлын биет
агуулга нь object storage (эсвэл локал файлын систем)-д хадгалагдаж,
thesis\_service нь зөвхөн pointer хадгалдаг.

\subsubsection{workflow\_service (Явцын Удирдлагын Сервис)}

Дипломын ажлын амьдралын мөчлөгийн статус шилжилтийг удирдана.
State Machine загварт суурилсан бөгөөд төлөвүүд нь:
\texttt{DRAFT → SUBMITTED → UNDER\_REVIEW → REVISION\_REQUESTED →
APPROVED → DEFENSE\_SCHEDULED → DEFENDED → ARCHIVED}.
Шилжилт бүр нь хугацааны тэмдэглэл, хариуцлагатай хэрэглэгч,
тайлбарын бичлэгийг агуулна.

\subsubsection{evaluation\_service (Үнэлгээний Сервис)}

Комиссын гишүүн бүрийн дүгнэлтийг хүлээн авч, нэгтгэлийн тооцоолол
хийн, эцсийн дүн гаргах. Дүн гаргах алгоритм нь тохируулгатай
(configurable) байхаар загварчлагдсан~--- жишилтийн жин, нийт онооны
хязгаар гэх мэт.

\subsubsection{committee\_service (Комиссын Удирдлагын Сервис)}

Хэлэлцүүлгийн комиссын бүрэлдэхүүн тогтоох, хяналтын хуваарь
гаргах, комиссын гишүүдэд мэдэгдэл илгээх. Энэ сервис нь
\texttt{user\_service}-аас гишүүдийн мэдээлэл авч,
\texttt{notification\_service}-тэй хамтарч мэдэгдэл явуулна.

\subsubsection{API Gateway}

Нэгдсэн нэвтрэх цэг (Single Entry Point) үүрэг гүйцэтгэж, гадаад
хүсэлт бүрийг хүлээн авч, JWT шалгаж, тохирох upstream сервис рүү
дамжуулна (proxy/reverse proxy). Rate limiting, request logging,
CORS хяналт нь gateway давхаргад хэрэгжинэ тул backend сервис тус
бүр энэ логикийг давтан хэрэгжүүлэх шаардлагагүй болно.

\begin{table}[H]
    \centering
    \caption{Backend микросервисүүдийн порт болон хариуцлагын хуваарь}
    \label{tab:service_ports}
    \renewcommand{\arraystretch}{0.7}
    \begin{tabular}{llll}
        \toprule
        \textbf{Сервис} & \textbf{Порт} & \textbf{Технологи} & \textbf{Гол хариуцлага} \\
        \midrule
        auth-service         & 8887 & Spring Boot  & JWT үүсгэх, нэвтрэлт \\
        user\_service        & 8086 & Spring WebFlux & Хэрэглэгчийн профайл \\
        topic\_service       & 8081 & Spring WebFlux & Сэдвийн удирдлага \\
        thesis\_service      & 8083 & Spring WebFlux & Дипломын ажлын мета \\
        workflow\_service    & 8084 & Spring WebFlux & Статус шилжилт \\
        evaluation\_service  & 8085 & Spring WebFlux & Дүгнэлт, үнэлгээ \\
        committee\_service   & 8082 & Spring WebFlux & Комиссын удирдлага \\
        Tomcat (JSF+MFE)     & 8080 & Apache Tomcat  & Host, static file serve \\
        \bottomrule
    \end{tabular}
\end{table}

\subsection{Сервисүүдийн Хоорондын Харилцааны Загвар}

Микросервисүүдийн хоорондын харилцааны загварчлалд хоёр үндсэн хэв
загвар хэрэглэгдэв:

\textbf{Синхрон харилцаа (Synchronous Communication):} HTTP/REST
протоколоор шууд дуудлага. \texttt{workflow\_service} нь статус
шилжилт хийхдээ \texttt{thesis\_service}-г шууд дуудаж дипломын
ажлын бичлэгийг шинэчлэнэ. Ийм шууд харилцааны тохиолдолд Circuit
Breaker загварыг (Nygard, 2007) хэрэглэж, upstream сервисийн
доройтол нь каскад байдлаар бусдад нэвтрэхээс хамгаалдаг.

\textbf{Асинхрон харилцаа (Asynchronous Communication):} Мэдэгдэл
(notification) явуулах, аудит бичлэг хийх гэх мэт цаг хугацааны
хамаарал бага үйлдлүүд нь Event-Driven загвараар хэрэгжих ёстой.
Ирээдүйн хөгжлийн шатанд Apache Kafka Message Broker нэвтрүүлснээр
\texttt{workflow\_service}-ийн статус шилжилтийн event нь
\texttt{notification\_service}-т асинхрон байдлаар хүрэх болно.

\begin{figure}[H]
    \centering
    \includegraphics[width=14cm]{fig_service_communication.png}
    \caption{Зураг 3.5. Микросервисүүдийн харилцааны диаграм. Хатуу шугам нь синхрон REST дуудлага, тасархай шугам нь асинхрон event-г тэмдэглэнэ. Эх сурвалж: зохиогчийн судалгаа, 2026.}
    \label{fig:service_communication}
\end{figure}

% ─────────────────────────────────────────────────────────────────────
\section{Өгөгдлийн Давхаргын Загварчлал}
% ─────────────────────────────────────────────────────────────────────

\subsection{Database Per Service загварын Сонголт}

Микросервисийн архитектурт өгөгдлийн давхаргыг хэрхэн хуваах вэ
гэсэн асуудал нь хамгийн нарийн төвөгтэй архитектурын шийдвэрүүдийн
нэг юм. Хоёр үндсэн сонголт байдаг:

\begin{enumerate}
  \item \textbf{Хуваалцсан өгөгдлийн сан (Shared Database):} Бүх
    сервис нэг өгөгдлийн санг хуваалцна. Хэрэгжүүлэлт энгийн боловч
    сервисүүдийн хооронд нягт холбоос (tight coupling) үүсдэг~---
    нэг сервисийн schema өөрчлөлт нь бусдад нөлөөлнө.
  \item \textbf{Сервис тус бүрийн өгөгдлийн сан (Database Per Service):}
    Сервис бүр өөрийн тусдаа schema эзэмшинэ. Сервис бүр бие даан
    deploy хийгдэх, технологийн сонголтоо чөлөөтэй хийх боломжтой
    боловч Cross-service join, distributed transaction нарийн болно.
\end{enumerate}

Энэхүү судалгаанд прагматик шийдлийг сонгосон: бүх сервис нэг
PostgreSQL instance-д ажилладаг боловч \textbf{schema-level тусгаарлалт}
хийгдсэн~--- сервис бүр өөрийн тусдаа schema-д бичиж, cross-schema
foreign key ашиглахгүй. Энэ нь ``Database Per Service'' загварын бие
даасан байдлын ашиг тусыг хэсэгчлэн хадгалахын зэрэгцээ нэг instance-ийн
хялбар засвар үйлчилгээний давуу талыг авч үлдэнэ.

\begin{figure}[H]
    \centering
    \includegraphics[width=14cm]{fig_database_schema_isolation.png}
    \caption{Зураг 3.6. Schema-level тусгаарлалтын загвар. Нэг PostgreSQL instance дотор сервис бүр өөрийн schema эзэмших бөгөөд хоорондын хамаарал нь application давхаргад шийдэгдэнэ. Эх сурвалж: зохиогчийн судалгаа, 2026.}
    \label{fig:db_schema_isolation}
\end{figure}

\subsection{Домейн Загварын Тодорхойлолт}

Системийн гол домейн объектуудын хоорондын харилцааг Entity-Relationship
(ER) загварт тодорхойлов. Хэдийгээр Cross-service join ашиглахгүй ч,
бизнесийн логикийн хэрэгцээгээр тодорхой домейн хооронд нягт харилцаа
оршино. Энэ харилцааг application давхаргад API дуудлагаар шийдвэрлэнэ.

Гол домейн объектууд болон тэдгээрийн атрибутуудыг дараах байдлаар
тодорхойлов:

\begin{itemize}
  \item \textbf{User} (user\_service): \texttt{id, sisiId, firstName,
    lastName, email, role, departmentId}
  \item \textbf{Department} (user\_service): \texttt{id, name, code,
    facultyId}
  \item \textbf{Topic} (topic\_service): \texttt{id, title, description,
    teacherId, maxStudents, status, academicYear}
  \item \textbf{Thesis} (thesis\_service): \texttt{id, title, studentId,
    teacherId, topicId, fileUrl, version, submittedAt}
  \item \textbf{WorkflowState} (workflow\_service): \texttt{id, thesisId,
    fromStatus, toStatus, actorId, comment, transitionedAt}
  \item \textbf{Committee} (committee\_service): \texttt{id, name,
    defenseDate, venueId, chairId}
  \item \textbf{CommitteeMember} (committee\_service): \texttt{id,
    committeeId, userId, role}
  \item \textbf{Evaluation} (evaluation\_service): \texttt{id, thesisId,
    committeeId, evaluatorId, score, criteria, submittedAt}
\end{itemize}

\begin{figure}[H]
    \centering
    \includegraphics[width=14cm]{fig_domain_er_diagram.png}
    \caption{Зураг 3.7. Системийн домейн загварын ER диаграм. Хатуу шугам нь нэг schema дотрох foreign key, тасархай шугам нь application давхаргын logical харилцааг заана. Эх сурвалж: зохиогчийн судалгаа, 2026.}
    \label{fig:domain_er}
\end{figure}

\subsection{Мэдээллийн Нийцэл ба Eventual Consistency}

Cross-service мэдээллийн нийцэлд ``Eventual Consistency'' (Аажмын
нийцэл) зарчмыг хэрэглэнэ. Distributed Transaction (2-Phase Commit)
нь микросервисийн орчинд хэдийгээр боломжтой боловч хүнд хэцүү
тогтолцоо шаарддаг тул Saga Pattern (Garcia-Molina \& Salem, 1987)
болон Outbox Pattern-ийг илүүд үзнэ.

Жишээ нь: оюутан дипломын ажлаа \texttt{SUBMITTED} төлөвт шилжүүлэхэд
\texttt{workflow\_service} нь шилжилтийн бичлэгийг хадгалж, дараа нь
\texttt{thesis\_service}-г шинэчлэх хүсэлт илгээнэ. Хэрэв
\texttt{thesis\_service} хариу өгөхгүй бол \texttt{workflow\_service}
нь retry логикоор дахин оролдох бөгөөд эцэст нь нийцэл тогтоно~---
энэ нь ``Eventual'' буюу аажмын нийцэл юм.

% ─────────────────────────────────────────────────────────────────────
\section{Нэгдсэн Архитектурын Загварын Нэгтгэл}
% ─────────────────────────────────────────────────────────────────────

Дээр тодорхойлсон гурван давхаргын загварчлал нэгдсэн байдлаар ажиллахад
дараах нэгдсэн зарчмуудыг дагана:

\begin{enumerate}
  \item \textbf{Нэг хариуцлага:} MFE модуль болон backend сервис тус
    бүр нэг бизнесийн чадварыг хариуцна~--- хариуцлагыг ``кэйс''
    дагуу биш, ``домейн'' дагуу хуваана.
  \item \textbf{API гэрээгээр харилцах:} Сервис болон MFE хоорондын
    харилцааны цорын ганц хэлбэр нь баримтжуулсан API гэрээ (OpenAPI
    Specification) юм~--- шууд function call, shared memory гэж байхгүй.
  \item \textbf{Бие даасан байршуулалт:} Сервис болон MFE тус бүр
    бусдыг зогсоолгүй шинэчлэгдэх чадвартай байна.
  \item \textbf{Барилгын бие даасан байдал:} Хөгжүүлэгч баг тус бүр
    өөрийн домейнийг бие даан эзэмшиж, технологийн сонголт дотор
    ухаалаг шийдэл гаргах эрхтэй.
\end{enumerate}

\begin{figure}[H]
    \centering
    \includegraphics[width=14cm]{fig_full_architecture_overview.png}
    \caption{Зураг 3.8. Системийн бүрэн архитектурын нэгдсэн бүдүүвч. Гурван давхарга (Frontend MFE, API Gateway, Backend Сервис) ба өгөгдлийн давхарга нэгдсэн байдлаар харагдана. Эх сурвалж: зохиогчийн судалгаа, 2026.}
    \label{fig:full_architecture}
\end{figure}

3-р бүлгийн загварчлалын шийдвэрүүд нь 4-р бүлгийн практик хэрэгжүүлэлтийн
суурь болно. Дараах бүлэгт ``юуг яагаад ийм загварчилсан'' гэдгийг
``хэрхэн бодитоор кодлон хэрэгжүүлсэн'' гэдэгтэй холбон тайлбарлана.
